\section{Cote 66: numerical, homological, and rational equivalence}
\label{sec:part-b}

We use rational coefficients unless integral coefficients are explicitly specified.
For a smooth projective variety $X/k$ of dimension $d$, a cycle
$z\in\mathrm{CH}^r(X)_{\Q}$ is numerically trivial if
$\deg(z\cdot w)=0$ for every $w\in\mathrm{CH}^{d-r}(X)_{\Q}$.
Its $\ell$-adic homological class means its class in
$H^{2r}_{\mathrm{\acute et}}(X_{\bar k},\Q_\ell(r))$, where
$\ell\ne\operatorname{char}k$. Thus we use geometric cohomology, not
Galois invariants or cohomology of $\operatorname{Spec}k$.
The official Montpellier scan, cote 66/004, asks both about
numerical versus $\tau$-equivalence for curves on a threefold and about
numerical versus homological equivalence for middle-dimensional cycles on
a fourfold. These are different questions.

\subsection{B.1(a) and L1: the characteristic-zero assertion}

\begin{theorem}[Proved; L1 over the complex numbers]\label{thm:B1a-betti}
Let $X$ be a smooth connected projective complex variety of dimension four.
Two codimension-two cycles with rational coefficients on $X$ are numerically
equivalent if and only if their rational Betti classes agree.
\end{theorem}
\begin{proof}
We give the algebraic Lefschetz decomposition explicitly, since this is the
step needed to apply Hodge--Riemann to numerical equivalence. Write
$h=c_1(H)\in H^2(X,\Q)$ for an ample line bundle, and
$A^2\subset H^4(X,\Q)$ for the image of the cycle-class map. Put
$V=\int_Xh^4>0$. For $\alpha\in A^2$ set
\[
 a=\frac{\int_X\alpha h^2}{V},\qquad
 b=(L^2:H^2(X,\Q)\longrightarrow H^6(X,\Q))^{-1}
       (L\alpha-a h^3),\qquad
 \gamma=\alpha-a h^2-hb,
\]
where $L$ denotes cup product with $h$. The inverse exists by hard Lefschetz.
It preserves rationality and shifts Hodge bidegree by $(-2,-2)$.
Consequently $b$ is a rational $(1,1)$-class. The Lefschetz $(1,1)$ theorem
makes $b$ a rational divisor class. In particular $\gamma$ is algebraic.
Moreover
\[
 L^3b=L^2\alpha-a h^4=0,\qquad L\gamma=0.
\]
The first equality uses $H^8(X,\Q)=\Q h^4$ and the definition of $a$.
Thus
\[
 A^2=\Q h^2\ \oplus\ h\operatorname{NS}(X)_{\Q,\mathrm{prim}}
          \ \oplus\ A^2_{\mathrm{prim}},
\]
where the primitive conditions are respectively $L^3b=0$ and $L\gamma=0$.
This decomposition is direct by the cohomological Lefschetz decomposition.

Let $q(u,v)=\int_Xuv$. The three displayed summands are mutually
orthogonal: $q(h^2,hb)=\int_Xh^3b=0$;
$q(h^2,\gamma)=0$ and $q(hb,\gamma)=0$ because $L\gamma=0$.
The form is positive definite on $\Q h^2$, negative definite on
$h\operatorname{NS}(X)_{\R,\mathrm{prim}}$, and positive definite on
$A^2_{\mathrm{prim}}\otimes\R$. These are precisely the signs in the
Hodge--Riemann relations for primitive $(1,1)$ and $(2,2)$ classes.
We use hard Lefschetz and Hodge--Riemann in the forms recorded in
\cite[\S2.2, Theorems 2.10--2.12]{BVoisin2010}; the divisor assertion is
the Lefschetz $(1,1)$ theorem recalled in its \S2.1.

If $\alpha$ is numerically trivial, it is orthogonal to every element of
$A^2$. We may therefore test against $a h^2$, $hb$, and $\gamma$, because
each is algebraic. Orthogonality of the decomposition and definiteness on
each summand give $a=0$, $b=0$, and $\gamma=0$. Hence $\alpha=0$.
Conversely a zero cohomology class has zero intersection number with every
complementary algebraic cycle by compatibility of cycle classes with
products and the trace map. This proves the assertion.
\end{proof}

The theorem is classical: Kleiman explicitly records the dimension-at-most-four
consequence in \cite[\S3.10, following Propositions 3.6 and 3.8 and Corollary
3.9]{BKleiman1968}. We claim no new result. The proof also explains why the
algebraicity of the Lefschetz components cannot simply be omitted.

\begin{lemma}[Proved; numerical triviality and field extension]\label{lem:B-basechange-num}
Let $X/k$ be smooth projective, $K/k$ a field extension, and
$z\in\mathrm{CH}^r(X)_{\Q}$. Then $z$ is numerically trivial if and only if
$z_K$ is numerically trivial.
\end{lemma}
\begin{proof}
If $z_K$ is trivial, test against pullbacks of cycles on $X$ and use invariance
of the degree of a zero-cycle under extension of its ground field.
For the converse let $w$ be a complementary integral cycle on $X_K$.
It is defined over a finitely generated intermediate field $k\subset K_0\subset K$.
Choose an integral finite-type $k$-scheme $T$ with function field $K_0$,
and spread $w$ to a cycle on $X\times_k T$. After shrinking $T$ its
components are flat over $T$. Intersection numbers with the fixed cycle
$z$ are constant under this specialization: intersection with a class on
the smooth factor, followed by proper pushforward to $T$, commutes with
specialization to its fibres. This is the conservation-of-number property
of intersection products. In this application one can compute it by a
finite locally free resolution of the fixed cycle on $X$: the Euler
characteristic of its derived tensor product with the flat family is
locally constant, and for complementary dimensions it is exactly the
intersection number. Thus at a closed point $t\in T$ we have
\[
 \deg_{K_0}(z_{K_0}w)=\deg_{k(t)}(z_{k(t)}w_t).
\]
The residue field $k(t)$ is finite over $k$. For
$p:X_{k(t)}\to X$, the projection formula gives
\[
 [k(t):k]\deg_{k(t)}(z_{k(t)}w_t)
     =\deg_k(z\,p_*w_t)=0.
\]
Division is legitimate in $\Q$. This proves the converse, and extension
from $K_0$ to $K$ again preserves the degree. The spreading operation and
each of the comparison maps used here preserve the intersection number
in the direction required; we do not assume that every cycle over $K$
descends to $k$.
\end{proof}

\begin{theorem}[Proved; B.1(a) and L1 over every characteristic-zero field]\label{thm:B1a-etale}
Let $k$ be any field of characteristic zero, $X/k$ smooth projective of
pure dimension four, and $\ell$ any prime. Numerical and geometric
$\ell$-adic homological equivalence coincide for codimension-two cycles
with rational coefficients on $X$.
\end{theorem}
\begin{proof}
The implication from homological to numerical equivalence follows from
products and traces. Conversely, descend $X$, a numerically trivial cycle
$z$, and an ample line bundle to a field $k_0\subset k$ finitely generated
over $\Q$. Denote the resulting objects by $X_0,z_0,H_0$.
Testing cycles defined over $k_0$ after extension to $k$ shows that $z_0$
is numerically trivial. Choose an embedding $k_0\hookrightarrow\C$.
Lemma~\ref{lem:B-basechange-num} shows that $z_{0,\C}$ is numerically
trivial, so Theorem~\ref{thm:B1a-betti}, applied component by component,
makes its Betti class zero. The Betti--\'etale comparison theorem
identifies its $\Q_\ell$ class with zero. Invariance of geometric
\'etale cohomology under extension of algebraically closed fields then
gives $\operatorname{cl}_\ell(z)=0$.
Here the comparison maps are the maps induced by
$X_{0,\C}\to X_{0,\bar k_0}$ and by
$X_{\bar k}\to X_{0,\bar k_0}$, with algebraic closures chosen compatibly;
they are isomorphisms and commute with cycle classes.
The comparison input is \cite[Expos\'e XVI, Th\'eor\`eme 4.1]{BSGA4XVI};
geometric extension invariance is its Corollary 1.6. Apply these first
to $\Z/\ell^n$ and then pass to inverse limits and invert $\ell$.
\end{proof}

\subsection{B.1(b): what Ancona proves, and the additional 2026 input}

\begin{proposition}[Proved; the precise positivity criterion]\label{prop:B1b-criterion}
Fix a Weil cohomology on smooth projective varieties over a field $k$.
Let $X/k$ be a smooth projective fourfold and $h$ an ample divisor class.
Suppose that its algebraic classes in $H^4$ admit the orthogonal
decomposition
\[
 A^2=\Q h^2\oplus hA^1_{\mathrm{prim}}\oplus A^2_{\mathrm{prim}},
\]
that the intersection form is negative definite on
$hA^1_{\mathrm{prim}}$ and positive definite on
$A^2_{\mathrm{prim}}$. Positivity here is on the $\Q$-vector spaces of
\emph{homological} algebraic classes, before dividing by numerical
equivalence. Then numerical and homological equivalence coincide in
codimension two. Conversely, if the decomposition exists, numerical
primitive positivity together with numerical--homological equality
implies this criterion.
\end{proposition}
\begin{proof}
The intersection form is positive on $\Q h^2$. A vector in its radical
on $A^2$ has, by orthogonality and definiteness, zero component in each
summand. The radical is exactly the space of numerical-zero homological
classes, since $X$ has dimension four. The converse follows by identifying
the numerical and homological spaces under the assumed equality.
The classical implication is also
\cite[Proposition 5-1, p.~17]{BKleiman1994}; its preceding definition
of $\operatorname{Hdg}(X)$ on pp.~16--17 is explicitly positivity
on algebraic cohomology classes. This is the positivity used here.
\end{proof}

\begin{proposition}[Proved; limitation of the proposed Ancona-only inference]\label{prop:B1b-gap}
The Lefschetz standard conjecture and positivity of the primitive form
\emph{on numerical classes} do not, by the preceding linear-algebra
argument alone, prove numerical--homological equality. Ancona's 2021
theorem supplies this numerical positivity for abelian fourfolds; its
stated equality theorem covers infinitely many primes $\ell$, not every
$\ell$.
\end{proposition}
\begin{proof}
Let $V=\Q e\oplus\Q r$ and let $q(e,e)=1$, $q(e,r)=q(r,r)=0$.
On $V/\operatorname{rad}(q)$ the induced form is positive definite,
although the radical of $V$ is nonzero. Adding independent positive and
negative definite Lefschetz summands changes nothing about this defect.
Thus positivity on the quotient cannot eliminate the radical before
quotienting. This is an explicit counterexample to the proposed
linear-algebra inference, not a counterexample to the standard conjecture.
Ancona defines his Hodge standard conjecture on numerical primitive
classes in \cite[\S3, Conjecture 3.10]{BAncona2021}; Theorem 3.18 proves it
for abelian fourfolds. Proposition 3.11 identifies the homological radical,
and Corollary 3.12 gives only semipositivity on the homological space.
Theorem 1.6, equivalently Theorem 3.20, gives the infinitely-many-primes
conclusion. These distinctions also appear explicitly in
\cite[Remark 1.3]{BBroe2026}.
\end{proof}

\begin{theorem}[Known with checked deduction from a 2026 preprint; B.1(b)]\label{thm:B1b-abelian}
Let $A/k$ be an abelian fourfold over any field and let
$\ell\ne\operatorname{char}k$. Numerical and $\ell$-adic homological
equivalence coincide for all cycles on $A$ with rational coefficients,
in particular for its two-dimensional cycles.
\end{theorem}
\begin{proof}
The additional input available as of 19 November 2026 is
\cite[Theorem 1.1]{BBroe2026}, which proves the Tate conjecture for
abelian fourfolds over finite fields. We explicitly use this opened
preprint as a known theorem; we do not claim an independent verification
of its full proof. Broe's Corollary 8.1(i), obtained from
\cite[Th\'eor\`eme 1, p.~978]{BKahn2003}, states that rational and numerical
equivalence coincide with rational coefficients over a finite field.
Since rational equivalence implies $\ell$-adic homological equivalence,
which implies numerical equivalence, this proves the required equality
on every finite-field fibre, for every invertible $\ell$.

To explain the extension step, spread an abelian fourfold and any fixed
cycle to an abelian scheme over an integral finite-type scheme over
$\Z[1/\ell]$, after passing to a finitely generated field of definition
and shrinking. A closed fibre $A_0$ has finite residue field. The
algebraicity of Lefschetz projectors for abelian varieties is
\cite[Theorem 2A11]{BKleiman1968}, also
\cite[Proposition 4-3(2)(c), p.~16]{BKleiman1994}.
Under these projectors, smooth proper
base change identifies the generic geometric cohomology with the special
geometric cohomology and takes algebraic primitive classes into algebraic
primitive classes, preserving the intersection form.
On $A_0$ this form on homological primitive classes is positive definite:
finite-field numerical--homological equality just proved identifies it
with the numerical form, and Ancona's Theorem 3.18 makes that form
positive definite. A subspace inherits definiteness, so
Proposition~\ref{prop:B1b-criterion} proves the equality on the generic
fibre. The divisor and one-cycle pieces follow from the divisor case and
algebraic Lefschetz. This is exactly the proof of Ancona's
Corollary 3.19, now with its missing finite-field hypothesis supplied for
every $\ell$. Arbitrary field extensions are handled by
Lemma~\ref{lem:B-basechange-num} and geometric base-change invariance.
The resulting statement is also recorded explicitly as
\cite[Corollaries 1.2 and 8.2]{BBroe2026}.
\end{proof}

\begin{proposition}[Proved under stated hypotheses; scope of the same argument]\label{prop:B1b-scope}
The positivity-and-specialization argument applies to every smooth
projective fourfold $X/k$ whose algebraic Lefschetz decomposition exists
and which has a smooth projective specialization $X_0$ for which
numerical--$\ell$-adic equality and numerical primitive positivity both
hold. It also proves numerical--$\ell$-adic equality on every smooth
projective fourfold dominated by a smooth projective fourfold $Y$ having
that equality, by a generically finite surjective morphism
$f:Y\to X$. In particular it applies to fourfolds dominated in this
sense by an abelian fourfold.
\end{proposition}
\begin{proof}
For specialization, the primitive algebraic homological space of $X$
injects into that of $X_0$ by smooth proper base change. The fibre form is
positive definite by the two hypotheses; its restriction is positive
definite. The algebraic decomposition and the divisor index theorem give
Proposition~\ref{prop:B1b-criterion}. These are explicit sufficient
hypotheses, not an assertion that all fourfolds satisfy them.

For the second assertion, if $z$ is numerically trivial on $X$, then
$f^*z$ is numerically trivial on $Y$, because for every complementary
cycle $w$ on $Y$ the projection formula gives
$\deg_Y(f^*z\cdot w)=\deg_X(z\cdot f_*w)=0$.
The pullback is defined on Chow groups because both varieties are smooth.
The assumed equality for $Y$ gives $f^*\operatorname{cl}_\ell(z)=0$.
The cohomological push--pull identity
$f_*f^*=\deg(f)\operatorname{id}$ then gives
$\operatorname{cl}_\ell(z)=0$, since the nonzero integer $\deg(f)$ is
invertible in $\Q_\ell$. The reverse implication always holds.
\end{proof}

\subsubsection{Further explicit fourfolds and the exact scope of the method}

Write $D^2_\ell(X)$ for numerical--$\ell$-adic homological equality
with rational coefficients in codimension two. The following transfer
rules give concrete further members of the class, without using the
desired equality as a hypothesis about the target.

\begin{proposition}[Proved; algebraic-correspondence transfer]
\label{B1b-correspondence}
Let $X,Y_1,\ldots,Y_m$ be smooth projective fourfolds over $k$, and
suppose $D^2_\ell(Y_i)$ holds. Suppose algebraic correspondences
$a_i\in\operatorname{CH}^4(X\times Y_i)_\Q$ and
$b_i\in\operatorname{CH}^4(Y_i\times X)_\Q$ satisfy
\[
 \sum_i(b_i)_*(a_i)_*=\operatorname{id}
       \quad\hbox{on }H^4(X_{\bar k},\Q_\ell(2)).
\]
Then $D^2_\ell(X)$ holds. In particular this applies to a generically
finite surjective morphism from any $Y_i$ to $X$.
\end{proposition}
\begin{proof}
If $z$ is numerically trivial on $X$, then $(a_i)_*z$ is numerically
trivial on $Y_i$: its intersection with a complementary cycle $w$
equals the intersection of $z$ with $({}^ta_i)_*w$, by the projection
formula. The hypothesis on $Y_i$ therefore gives
$(a_i)_*\operatorname{cl}_\ell(z)=0$. The displayed identity implies
$\operatorname{cl}_\ell(z)=0$. For $f:Y\to X$ of degree $d>0$, take
$a={}^t\Gamma_f$ and $b=d^{-1}\Gamma_f$ and use $f_*f^*=d$.
This includes inseparable maps: only the nonzero degree in $\Q_\ell$
is used. Homological triviality always implies numerical triviality.
\end{proof}

\begin{lemma}[Proved; smooth blow-ups of fourfolds]
\label{B1b-blowup}
Let $X$ be a smooth projective fourfold over $k$, let $Z\subset X$ be
a smooth closed center of positive codimension on every component,
and let $f:W=\operatorname{Bl}_ZX\to X$.
Then $D^2_\ell(W)$ holds if and only if $D^2_\ell(X)$ holds.
\end{lemma}
\begin{proof}
All cohomology below is geometric; the cycles and operations remain
defined over $k$, so the hypothesis $D^2_\ell(X)$ is used over $k$.
A Cartier center gives an isomorphism, so suppose its codimension
is $c\ge2$, component by component. Let $j:E\hookrightarrow W$
be the exceptional divisor, $q:E=\mathbf P(N_{Z/X})\to Z$, and
$\xi=c_1(\mathcal O_E(1))$. The blow-up decomposition
\cite[Lemma~2.1, p.~103]{ShiodaKatsura1979C} is
\[
 H^4(W,\Q_\ell(2))=
 f^*H^4(X,\Q_\ell(2))\oplus
 \bigoplus_{a=1}^{c-1}
 j_*\bigl(\xi^{a-1}q^*H^{4-2a}(Z,\Q_\ell(2-a))\bigr).
\]
All maps shown are induced by algebraic cycle operations.
Every smooth projective variety of dimension at most two has
numerical--homological equality in every codimension: for divisors
use Matsusaka's numerical--$\tau$ equality, recorded in
\cite[p.~361]{BKleiman1968}, and the vanishing of the rational first
Chern class of an algebraically trivial line bundle; zero-cycles
are measured by degree and codimension-zero classes by components.
This applies to $Z$.

Suppose $D^2_\ell(X)$ and let $z$ be numerically trivial on $W$.
Write its cohomology class in the displayed decomposition as
\[
 \alpha=f^*\alpha_0+
      \sum_{a=1}^{c-1}j_*(\xi^{a-1}q^*\beta_a).
\]
Since $f_*\alpha=\alpha_0$ and $f_*z$ is numerically trivial,
$\alpha_0=0$. The self-intersection identity $j^*j_*=-\xi$ gives,
for $b=c-1,c-2,\ldots,1$,
\[
 -q_*\bigl(\xi^{c-1-b}j^*\alpha\bigr)
 =\beta_b+\sum_{a>b}
       q_*(\xi^{c-1+a-b})\,\beta_a.
\]
The left side is the class of a numerically trivial cycle on $Z$:
proper pushforward, pullback between smooth varieties and multiplication
by $\xi$ preserve numerical triviality by the projection formula.
It is therefore zero. Descending induction gives every $\beta_b=0$;
groups of negative degree are simply zero. Hence $\alpha=0$.
Conversely apply Proposition~\ref{B1b-correspondence} to the degree-one
morphism $f$. No factorization theorem for arbitrary birational maps
in positive characteristic is used here.
\end{proof}

\begin{theorem}[Proved with checked source construction; Fermat cubic fourfolds]
\label{B1b-fermat}
Let $\operatorname{char}k\ne3$ and let
\[
 F=\{z_0^3+\cdots+z_5^3=0\}\subset\mathbf P^5_k.
\]
For every $\ell\ne\operatorname{char}k$, numerical and $\ell$-adic
homological equivalence agree for all cycles on $F$ with rational
coefficients. The same holds for any smooth projective $k$-form of $F$.
Together with Ancona's numerical Hodge theorem for $F$, this places
these fourfolds in the homological-positivity class of
Proposition~\ref{prop:B1b-criterion}.
\end{theorem}
\begin{proof}
We first work over an algebraic closure. Let
$S=\{x_0^3+x_1^3+x_2^3+x_3^3=0\}\subset\mathbf P^3$ and
$E=S\cap\{x_3=0\}$. Both are smooth because $3\ne0$ in $k$.
The cohomology of $S$ is algebraic; here is an explicit verification
valid also in characteristic two. Fix a primitive cube root $\zeta$.
For $a,b\in\Z/3$ put
\[
 L^0_{ab}: x_0+\zeta^a x_1=x_2+\zeta^b x_3=0,
 \qquad
 L^1_{ab}: x_0+\zeta^a x_2=x_1+\zeta^b x_3=0.
\]
These are lines on $S$. Within either family, two distinct lines meet
exactly when their first or their second indices agree; a line
$L^0_{ab}$ meets $L^1_{cd}$ exactly when $a+d=b+c$ in $\Z/3$.
These assertions follow by solving the four linear equations. They
hold in every characteristic different from three: the residual
equation is a difference of two distinct cube roots, which is nonzero.
Distinct meeting lines have intersection number one, and adjunction
gives self-intersection $-1$.

In the order
$L^0_{00},L^0_{01},L^0_{02},L^0_{10},L^0_{11},L^1_{00},L^1_{01}$,
the intersection matrix is
\[
 G=\begin{pmatrix}
 -1&1&1&1&0&1&0\\
 1&-1&1&0&1&0&1\\
 1&1&-1&0&0&0&0\\
 1&0&0&-1&1&0&0\\
 0&1&0&1&-1&1&0\\
 1&0&0&0&1&-1&1\\
 0&1&0&0&0&1&-1
 \end{pmatrix},\qquad \det G=1.
\]
The attached exact program checks this determinant and all 351 pairs
of the 27 coordinate-pairing lines over $\mathbf F_4$.
The preceding cube-root calculation proves the displayed incidence
matrix in every characteristic different from three; the finite program
is not used to infer this characteristic-independent assertion.
Over $\C$, weak Lefschetz gives
$b_1=b_3=0$, and
$c(TS)=(1+H)^4/(1+3H)$ gives
$\int_Sc_2(TS)=\int_S3H^2=9$, hence $b_2=7$.
For each fixed $\ell$, the smooth proper Fermat family over
$\Z[1/(3\ell)]$ and smooth proper base change give the same Betti
numbers in every characteristic considered.
Thus the seven independent line classes span $H^2(S,\Q_\ell(1))$;
$H^0$ and $H^4$ are also algebraic and odd cohomology is zero.
K\"unneth now shows that all the cohomology of $S\times S$ is
algebraic. Poincar\'e duality then gives $D^2_\ell(S\times S)$:
an algebraic class pairing to zero with every algebraic cycle pairs
to zero with the whole complementary cohomology space.

Shioda--Katsura's explicit construction
\cite[Lemmas~1.1--1.2, pp.~98--99; Theorem~1.7, pp.~101--102]
{ShiodaKatsura1979C}, with degree three and $r=s=2$, supplies a morphism
\[
 \psi:\operatorname{Bl}_{E\times E}(S\times S)\longrightarrow F
\]
of generic degree three. Explicitly its rational map before the blow-up is
\[
 ([x],[y])\longmapsto
 [x_0y_3:x_1y_3:x_2y_3:-x_3y_0:-x_3y_1:-x_3y_2].
\]
The sum of the cubes is zero by the two equations of $S$. Its base
ideal is locally $(x_3,y_3)$ along $E\times E$, so blowing up this
smooth codimension-two center resolves the map. On the dense affine
chart $x_3=y_3=1$, an inverse to a general projective target point
requires a scalar $a$ satisfying
$a^3=-1/(z_0^3+z_1^3+z_2^3)$; its three distinct choices give generic
degree three. Properness makes the morphism surjective.
Lemma~\ref{B1b-blowup} proves $D^2_\ell$ for its source, and
Proposition~\ref{B1b-correspondence} proves it for $F$.

Outside middle degree, the even cohomology of a smooth hypersurface
$F$ is generated by powers of its hyperplane class and its odd
cohomology vanishes \cite[Lemma~2.3, p.~103]{ShiodaKatsura1979C}.
This proves equality for all codimensions. Geometric base change and
Lemma~\ref{lem:B-basechange-num} give the statement over $k$, including
forms of $F$.

Finally \cite[Proposition~A.9]{BAncona2021} proves numerical primitive
positivity for this exact Fermat fourfold in every characteristic
different from three. Its algebraic Lefschetz splitting is explicit:
$H^2=\Q_\ell h$, $H^6=\Q_\ell h^3$, $\deg h^4=3$, and
\[
 \alpha=\frac{\deg(\alpha h^2)}{3}h^2+
       \left(\alpha-\frac{\deg(\alpha h^2)}{3}h^2\right),
\]
with the second term primitive. Equality already proved identifies
this numerical primitive form with its homological form. It is
therefore positive definite, as required.
\end{proof}

Consequently an explicit covered class consists of abelian fourfolds,
Fermat cubic fourfolds and their geometric forms, together with all
fourfolds obtained by finitely many applications of
Proposition~\ref{B1b-correspondence} and Lemma~\ref{B1b-blowup}.
The seeds also furnish special fibres for
Proposition~\ref{prop:B1b-scope}, whenever its actual smooth-family
and algebraic-Lefschetz hypotheses hold. An arbitrary fourfold is not
asserted to specialize to one of these seeds with its relevant cycle
data.

The exact domain of the \emph{componentwise definiteness argument}
is intrinsically specified by Proposition~\ref{prop:B1b-criterion}:
the Lefschetz projections must preserve algebraic classes and the
signed primitive intersection form on \emph{homological} algebraic
classes must be positive definite. These are independent cohomological
and positivity hypotheses; numerical--homological equality is the
conclusion, not an assumed property of the target. Under the additional
numerical-positivity hypothesis, the converse in that proposition
identifies the remaining issue precisely with the absence of a
homological radical. The conditional method criterion is complete.
The examples and transfer rules above certify substantial geometric
membership; they are not an assertion that this list exhausts all
geometric fourfolds in that intrinsic class. Ancona's
\cite[Remark~A.10]{BAncona2021} likewise distinguishes his rank-two
method's proved examples from a classification of all fourfolds.


\subsection{B.1(c): the first question of 66/004}

\begin{theorem}[Refuted; numerically trivial curves need not be $\tau$-trivial]\label{thm:B1c}
With the classical definition of $\tau$-triviality (some nonzero integral multiple is
algebraically equivalent to zero), the first assertion of question 24 is
false, even for smooth complex quintic threefolds. It is also false if the
doubtful glyph is read as algebraic triviality itself. The implication
from $\tau$-triviality to numerical triviality remains true.
\end{theorem}
\begin{proof}
We specify a generic counterexample completely. In $\mathbf P^4$ let
\[
 L_1=V(x_2,x_3,x_4),\qquad L_2=V(x_0,x_1,x_4).
\]
These are disjoint lines. Let $\mathcal M$ be the 114 degree-five
monomials vanishing on both lines. Take algebraically independent
indeterminates $t_m$ over $\Q$, indexed by $m\in\mathcal M$, put
$k=\overline{\Q(t_m:m\in\mathcal M)}$, and define
\[
 X=V\left(\sum_{m\in\mathcal M}t_m m\right)\subset\mathbf P^4_k,
 \qquad z=[L_1]-[L_2]\in\mathrm{CH}^2(X).
\]
This generic hypersurface is smooth. Indeed the same linear system has
the smooth member
\[
 F_0=x_0^4x_2+x_0x_2^4+x_1^4x_3+x_1x_3^4+x_4^5.
\]
For a pair $(a,b)$ the two partial derivatives of $a^4b+ab^4$ are
$b(4a^3+b^3)$ and $a(a^3+4b^3)$. They vanish simultaneously only at
$(0,0)$ in characteristic zero: if $ab\ne0$, the two linear equations
in $a^3,b^3$ have determinant $15\ne0$. The remaining derivative is
$5x_4^4$. Hence the five derivatives have no common projective zero.
Smoothness is open in the parameter space. A smooth hypersurface of
positive dimension in projective space is geometrically connected here,
and $X$ is projective of dimension three.

Weak Lefschetz and Poincar\'e duality give $H^4(X_\C,\Q)$ rank one for
any embedding of this characteristic-zero field of definition into $\C$.
Both lines have degree one, so their classes agree; consequently $z$ is
numerically trivial. Equivalently, weak Lefschetz gives
$\operatorname{Pic}(X)=\Z\mathcal O_X(1)$ and testing against its
generator gives zero.

The nontriviality input is \cite[\S III, assertion (3.3), p.~290]{BCeresaCollino1983}:
on a general quintic threefold no nonzero integral linear combination
of its 2875 lines is algebraically trivial. In particular
$n(L_1-L_2)$ is not algebraically trivial for any $n\ne0$ on a general
member of the ordered-skew-line incidence family. We justify passage
to our specified family. The space of ordered skew lines is a
homogeneous space for $\operatorname{PGL}_5$; the degree-five equations
vanishing on each pair form a vector bundle of rank 114. Its total
incidence space is therefore integral. The projection to quintic
hypersurfaces is dominant: a general quintic has skew lines, as recalled
at the start of \S III of that paper. Transitivity of
$\operatorname{PGL}_5$ identifies the fibres over all ordered skew pairs.
Thus a relation on the generic member with our fixed pair would, after
spreading its finite algebraic-equivalence data, hold on a nonempty
open subset of that incidence space. Its image contains an open subset
of the hypersurface parameter space. This contradicts (3.3), which
holds on a very general complex fibre. The data defining a relation
are algebraic and finite, so the same argument applies after algebraic
closure of the generic field. This proves that our $z$ is not
$\tau$-trivial. Choosing an embedding $k\hookrightarrow\C$ with generic
coefficients gives the asserted complex counterexample; algebraic
equivalence after this extension would likewise spread and give the
forbidden generic relation.

Finally an algebraically trivial cycle has zero intersection number
with every complementary cycle, since degree is constant in a connected
family. If $nz$ is algebraically trivial, then
$n\deg(zw)=0$ for every $w$, so $z$ is numerically trivial.
The program \path{code/part_b_finite_checks.py} verifies the monomial
count, the containment of both lines, and smoothness of $F_0$ over $\Q$.
It does not purport to decide algebraic equivalence computationally.
\end{proof}

\subsection{B.2: the square problem and the Albanese kernel}

On the official scan cote 66/005, question 25 first asks whether the
product of two algebraically trivial cycles is rationally trivial
(\emph{``pb du carr\'e''}); it then asks whether an Albanese-trivial
zero-cycle is rationally trivial. The marginal reference to Mumford
is adjacent to the surface zero-cycle question. A square on a surface
which can answer the first question must be the square of a divisor:
the square of a zero-cycle lies in $\mathrm{CH}^4$ and is zero for
dimension reasons.

\begin{lemma}[Proved from Mumford; moving pairs and two-forms]\label{lem:B2-mumford}
Let $S$ be a smooth projective complex surface, let
$\omega\in H^0(S,\Omega_S^2)$, and let $T$ be smooth and integral.
For a morphism $f:T\to\operatorname{Sym}^d S$, write $f^*\omega^{(d)}$
for Mumford's induced two-form. If $f,g:T\to\operatorname{Sym}^d S$
are families whose cycles are rationally equivalent at the generic
point after an algebraic extension of $\C(T)$, then
$f^*\omega^{(d)}=g^*\omega^{(d)}$.
\end{lemma}
\begin{proof}
Mumford constructs induced forms, including across diagonals of the
symmetric product, in \cite[\S1, pp.~197--199]{BMumford1969}; his Lemma 2
states their additivity under addition of cycles. His Lemma 3,
pp.~201--202, realizes a rational equivalence after a dominant change of
parameter by a family of rational curves in a symmetric product, after
adding the same auxiliary effective cycle to both endpoints.
Resolve the parameter space in characteristic zero and restrict to its
smooth dense open; call it $T'$, with dominant map $u:T'\to T$.
If $c$ is the auxiliary cycle and
$F:T'\times\mathbf P^1\to\operatorname{Sym}^{d+e}S$ is that family,
then
\[
 F|_0=f\circ u+c,\qquad F|_\infty=g\circ u+c.
\]
If the rational equivalence uses several such families, we apply the
following argument to each consecutive pair and add the equalities.
Every global two-form on $T'\times\mathbf P^1$ is pulled back from
$T'$, because $H^0(\mathbf P^1,\Omega^1)=
H^0(\mathbf P^1,\Omega^2)=0$. Therefore its restrictions at $0$ and
$\infty$ agree. Applying this to $F^*\omega^{(d+e)}$ and using
additivity cancels the form induced by $c$ and gives
$u^*f^*\omega^{(d)}=u^*g^*\omega^{(d)}$.
Pullback of differential forms by a dominant map of smooth
characteristic-zero varieties is injective at the generic point,
because the corresponding function-field extension is separable.
The desired equality follows. This is the argument of Mumford's
theorem, pp.~200 and 202--203, applied to two varying endpoints.
\end{proof}

\begin{theorem}[Refuted with complete counterexample; both questions 25]\label{thm:B2-square}
Let $E/\overline\Q$ be the smooth elliptic curve
$y^2=x^3-x$ with origin $o$, and let $S=E_\C\times E_\C$.
Choose $p,q\in E(\C)$ such that the point $(p,q)$ is generic over
$\overline\Q$; for example choose $x(p),x(q)$ algebraically independent
over $\overline\Q$ and then choose the square roots defining their
$y$-coordinates. Put
\[
 a=\operatorname{pr}_1^*([p]-[o]),\quad
 b=\operatorname{pr}_2^*([q]-[o]),\quad
 D=a+b\in\mathrm{CH}^1(S).
\]
Then $D$ is algebraically trivial, but $D^2$ is nonzero, and indeed
nontorsion, in $\mathrm{CH}^2(S)$. Moreover
\[
 z=[(p,q)]-[(p,o)]-[(o,q)]+[(o,o)]
\]
is Albanese-trivial and nontorsion. Thus each question has a negative
answer with integral coefficients and with rational coefficients.
\end{theorem}
\begin{proof}
The cubic defining $E$ has distinct roots in characteristic zero, so
its smooth projective completion is an elliptic curve. Divisors
$[p]-[o]$ and $[q]-[o]$ have degree zero and move in connected families
parametrized by $E$; they are algebraically trivial. Pullback and addition
preserve this property, hence $D$ is algebraically trivial.
Since each of $a,b$ is pulled back from a curve,
$a^2=b^2=0$ in the Chow ring of the surface. The fibre divisors from the
two factors meet transversally, so
\[
 ab=z,\qquad D^2=2ab=2z.
\]
The Albanese map of $E\times E$, based at $(o,o)$, is the identity map
to the abelian variety $E\times E$. On degree-zero zero-cycles it is
addition of the points with their coefficients. Therefore
\[
 \operatorname{alb}(z)=(p,q)-(p,o)-(o,q)+(o,o)=(o,o).
\]
These statements use the definitions of rational and algebraic
equivalence and of the Albanese homomorphism in
\cite[Lecture 0, pp.~3--5, and Lecture 1, Proposition 1.1(ii)--(iii),
pp.~9--10]{BBloch1980}; no
identification of the Albanese kernel with zero is being assumed.

We prove that $nz\ne0$ for every integer $n>0$. Consider the maps
from $S$ to $\operatorname{Sym}^2 S$
\[
 f(p',q')=[(p',q')]+[(o,o)],\qquad
 g(p',q')=[(p',o)]+[(o,q')].
\]
Let $\eta$ be a nonzero regular one-form on $E$ and let
$\omega=\operatorname{pr}_1^*\eta\wedge\operatorname{pr}_2^*\eta$.
This is a nonzero regular two-form on $S$.
The induced forms are
\[
 f^*\omega^{(2)}=\omega,\qquad g^*\omega^{(2)}=0:
\]
the constant summand contributes zero, and each summand of $g$ factors
through a curve. On multiplying both effective cycles by $n$,
additivity changes these forms to $n\omega$ and zero.

For each $n>0$, auxiliary degree $e\geq0$, and degree $d$ of a
rational curve in the projective variety $\operatorname{Sym}^{2n+e}S$,
form over $\overline\Q$ the finite type incidence scheme
\[
 W_{n,e,d}=\{(t,c,h):h(0)=nf(t)+c,\quad h(\infty)=ng(t)+c\},
\]
where $c\in\operatorname{Sym}^eS$ and $h:\mathbf P^1\to
\operatorname{Sym}^{2n+e}S$ has degree $d$.
Mumford's Lemma 3 shows that every rational equivalence at issue
occurs in one of these parameter schemes. If an irreducible component
of a $W_{n,e,d}$ had dense image in $S$, its endpoint family would
satisfy Lemma~\ref{lem:B2-mumford} over a dominant parameter space,
forcing $n\omega=0$, a contradiction. Each image therefore has proper
closed Zariski closure defined over $\overline\Q$. This argument uses
constructibility of each image, and does not need a separate assertion
that the image itself is closed. Our specified point $(p,q)$,
which is generic over $\overline\Q$, lies in none of them. Hence
$nz\not\sim_{\mathrm{rat}}0$ for every $n>0$, as required.
In particular $D^2=2z$ is nonzero even after tensoring with $\Q$.
\end{proof}

\begin{proposition}[Proved; exact positive replacement on this surface]\label{prop:B2-kernel}
For $S=E_\C\times E_\C$, the kernel of
$\operatorname{alb}:\mathrm{CH}_0(S)^0\to\operatorname{Alb}(S)(\C)$
on degree-zero classes is generated by the cycles
\[
 z(p,q)=([p]-[o])\times([q]-[o]),\qquad p,q\in E(\C).
\]
Consequently over $\Q$ this kernel is spanned by squares of
algebraically trivial divisors. Products in the counterexample are
Albanese-trivial, although they need not be rationally trivial.
\end{proposition}
\begin{proof}
For a degree-zero cycle $c=\sum_i n_i[(p_i,q_i)]$ the identity
\[
 c=\sum_i n_i z(p_i,q_i)
       +\sum_i n_i[(p_i,o)]+\sum_i n_i[(o,q_i)]
\]
holds because $\sum_i n_i=0$. If $\operatorname{alb}(c)=0$, then
$\sum_i n_i p_i=o$ and $\sum_i n_i q_i=o$ in $E$.
Abel's theorem for a smooth genus-one curve says that a degree-zero
divisor with this sum is principal; see \cite[Lecture 0, pp.~3--4;
Lecture 1, p.~11]{BBloch1980}. Thus both final sums are rationally
trivial on their copies of $E$, and also after pushforward to $S$.
Each $z(p,q)$ is itself Albanese-trivial by the calculation above, so
these generators give the whole kernel. Finally
$2z(p,q)=(a+b)^2$ with the corresponding pullback divisors $a,b$.
After division by two this proves the square-spanning assertion.
\end{proof}

The relevant Mumford paper was communicated on 28 November 1968 but
published in \emph{J.~Math.~Kyoto Univ.} 9(2) (1969), pp.~195--204.
We retain the journal year in the bibliography. The construction does
not assert that an arbitrary pair of nontorsion points has nonzero
exterior product; the proved genericity condition is part of the data.

\subsection{B.3, subject 21: the exact affine characterization}

\begin{theorem}[Proved from Gabber; universal degree-zero test]\label{thm:B3-21}
Let $A$ be a commutative ring, $I\subset A$ an ideal,
$X=\operatorname{Spec}A$, $Z=\operatorname{Spec}(A/I)$, and $i:Z\to X$.
The following statements are equivalent:
\begin{enumerate}[label=(\roman*)]
\item $(A,I)$ is a henselian pair.
\item For every torsion abelian \'etale sheaf $F$ on $X$, the restriction
$H^0(X,F)\to H^0(Z,i^*F)$ is an isomorphism.
\item For every such $F$ and every $q\ge0$, the restriction
$H^q(X,F)\to H^q(Z,i^*F)$ is an isomorphism.
\end{enumerate}
In particular $I$-adic completeness is sufficient, without a noetherian
hypothesis. The quantifier ``every sheaf'' cannot be replaced by ``one
fixed sheaf'' in the degree-zero test.
\end{theorem}
\begin{proof}
The implication (i)$\Rightarrow$(iii) is Gabber's affine analogue of
proper base change, with no restriction on the torsion primes: the
complete proof is opened in \cite[Theorem 59.82.7, tag 09ZI]{BStacks},
following Lemma 59.82.6 for degree zero. The implication
(iii)$\Rightarrow$(ii) is immediate.

Assume (ii). For every finite $A$-algebra $B$, let
$f:Y=\operatorname{Spec}B\to X$ and take
$F=f_*\underline{\Z/2\Z}$. This is a torsion abelian \'etale sheaf.
The degree-zero case of proper base change for the finite morphism $f$
identifies $i^*F$ with the direct image of
$\underline{\Z/2\Z}$ on
$Y\times_X Z=\operatorname{Spec}(B/IB)$.
Thus (ii) says that the natural map
\[
 H^0(Y,\underline{\Z/2\Z})\longrightarrow
 H^0(\operatorname{Spec}(B/IB),\underline{\Z/2\Z})
\]
is bijective. A section of this constant two-element sheaf is exactly
a clopen decomposition, hence exactly an idempotent in the ring.
Therefore idempotents of $B$ and $B/IB$ correspond bijectively for every
finite $A$-algebra $B$. By the henselian-pair criterion
\cite[Lemma 15.11.6, tag 09XI]{BStacks}, this is equivalent to (i).
That criterion includes the Jacobson-radical condition on $I$; it has
not been silently dropped.

An $I$-adically complete ring is henselian by
\cite[Lemma 15.11.4, tag 0ALJ]{BStacks}, so completeness gives the
stated consequence. For the last assertion let
$X=\mathbf G_{m,\C}$, $Z=\{1\}$, and $F=\underline{\Z/n\Z}$ with
$n>1$. Both spaces are connected, so their degree-zero groups are
$\Z/n\Z$ and restriction is an isomorphism. But the Kummer torsor
$u\mapsto u^n$ gives
$H^1(X,F)=\Z/n\Z$, whereas $H^1(Z,F)=0$.
Thus the fixed-sheaf interpretation is false.
\end{proof}

\begin{theorem}[Known with checked application; the distinct completion theorem]\label{thm:B3-fujiwara}
Let $(A,I)$ be a henselian pair with $I$ finitely generated, let
$\widehat A$ be its $I$-adic completion, and let
$U\subset\operatorname{Spec}A$ be an open subset containing
$\operatorname{Spec}A\setminus V(I)$.
Put $\widehat U=U\times_A\widehat A$ and let
$\pi:\widehat U\to U$ be the projection.
For every torsion abelian \'etale sheaf $F$ on $U$ and every $q\ge0$,
\[
 H^q(U,F)\ \xrightarrow{\ \sim\ }\
 H^q(\widehat U,\pi^*F).
\]
\end{theorem}
\begin{proof}
Apply \cite[Corollary 2.36, with Notation 2.33]{BHaineEtAl2024} to
the Eilenberg--Mac Lane sheaf of spaces $K(F,q)$ and take connected
components of the resulting equivalence on global sections.
This is the modern noetherian-free form of Fujiwara--Gabber rigidity.
In the opened paper, the classical source is identified as Fujiwara
1995, Corollary 6.6.4; the abelian theorem in the generality just stated
is also derived there from Bhatt--Mathew, Theorem 6.11.
For the basic case $U=\operatorname{Spec}A\setminus V(I)$,
the opened modern proof can be made explicit as follows.
Extend $F$ by zero to a torsion sheaf $G$ on $X=\operatorname{Spec}A$.
By \cite[Theorems 5.4 and 6.4, author pp.~38 and 48]{BBhattMathewArc},
torsion \'etale cohomology is a finitary arc-sheaf and the formal
glueing square
\[
\begin{array}{ccc}
R\Gamma(X,G)&\longrightarrow&R\Gamma(\operatorname{Spec}\widehat A,G)\\
\big\downarrow&&\big\downarrow\\
R\Gamma(U,F)&\longrightarrow&R\Gamma(\widehat U,\pi^*F)
\end{array}
\]
is homotopy cartesian. In this diagram all sheaves on the right are
pulled back from the corresponding sheaves on the left.
The upper arrow is an equivalence by Theorem~\ref{thm:B3-21},
applied to $(A,I)$ and $(\widehat A,I\widehat A)$, whose quotients
are canonically the same. Therefore the lower arrow is an
equivalence. This is exactly the derivation in
\cite[Corollary 6.10 and Theorem 6.11, author p.~50]{BBhattMathewArc}.
Its Remark 6.12 states explicitly that Fujiwara's original theorem
assumed $A$ noetherian; the general nonnoetherian result is due
to Gabber. The arbitrary containing open $U$ in our statement is
covered by the above direct application of Corollary 2.36.
Unlike Theorem~\ref{thm:B3-21}, this is a theorem about the complement
of a closed set and passage to completion, not restriction to the
closed fibre. The source's Remark 2.37 makes this distinction explicit.
\end{proof}

We read ``en dim 0'' in question 21 as cohomological degree zero in
Theorem~\ref{thm:B3-21}; if it means schemes of dimension zero, that is a
different assertion and is not used in the proof. Gabber's original
1994 article and Fujiwara's original 1995 article have bibliographic
records checked, but their original full texts have not yet been
retrieved in this run. The modern full proofs just cited have been
opened. Accordingly the particular demand to derive the conclusion
directly from both original articles remains a source-audit gate,
although the mathematical affine assertion above is proved.
The hypotheses of the two verified modern statements are distinct:
Theorem~\ref{thm:B3-21} allows an arbitrary ideal in an arbitrary
commutative ring, whereas Theorem~\ref{thm:B3-fujiwara} requires the
ideal to be finitely generated. Neither statement assumes
noetherianity or excludes torsion at residue characteristics.
We use the opened Stacks proof for the former unrestricted
generality, and do not infer the exact hypotheses of Gabber's
unopened original from its bibliographic abstract. The latter
theorem's original noetherian scope is explicitly recorded in the
opened Bhatt--Mathew Remark 6.12.

The literal named-source list also includes Raynaud's published thesis
\cite{BRaynaudThesis1970}, namely \emph{Faisceaux amples sur les
sch\'emas en groupes et les espaces homog\`enes}, LNM~119 (1970),
based on his Paris thesis of 1968. Its publisher contents and original
front matter have been opened, but its mathematical text has not.
It is distinct from the 1970 Picard article actually used below.
Lieberman's 1968 article likewise remains an original-source reading
requirement; the abelian Lefschetz result used here is instead checked
in the opened Kleiman originals. These source-access requirements do
not change the status of the deductions proved above.

\subsection{B.3: precise answers to the programme headings}

The remaining entries of cote 66 are mostly programme headings rather
than propositions with fixed hypotheses. We therefore give an actual
theorem, its category and coefficient conditions, and the deduction
that connects it to each heading. A restricted theorem below is not
silently promoted to a solution of a broader interpretation. The
original scan 66/002 has been opened: in particular subject 11 really
reads ``sch\'emas morcel\'es'', and subject 3 says ``partie
$p$-primaire'', not ``Swan conductor''. These distinctions matter.

\begin{theorem}[Known with checked deduction; subjects 1 and 8]\label{thm:B3-flat-crys}
For any scheme $X$, let
$a_X:\operatorname{Sh}((\operatorname{Sch}/X)_{\mathrm{fppf}})
\to\operatorname{Sh}(X_{\mathrm{et}})$ be the comparison morphism.
For every bounded below complex $K$ of abelian \'etale sheaves,
\[
 K\simeq Ra_{X,*}a_X^*K,
 \qquad
 R\Gamma(X_{\mathrm{et}},K)
 \simeq R\Gamma(X_{\mathrm{fppf}},a_X^*K).
\]
There is no invertibility condition on the torsion in this assertion.

For the crystalline addition to subject 1 and the connection addition
to subject 8, let $(A,I,\gamma)$ be a divided power ring with a prime
$p$ nilpotent in $A$. Put $S=\operatorname{Spec}A$, let $Y\to S$
be smooth, and put $X=Y\times_A A/I$. If $\mathcal F$ is a crystal
in quasi-coherent $\mathcal O_{X/S}$-modules, its value $\mathcal F_Y$
has a canonical integrable connection and
\[
 R\Gamma(\operatorname{Cris}(X/S),\mathcal F)
 \simeq R\Gamma(Y,\mathcal F_Y\otimes\Omega^\bullet_{Y/S}).
\]
\end{theorem}
\begin{proof}
The first unit isomorphism is
\cite[Lemma 59.100.7, in tag 0DDK]{BStacks}. Applying derived global
sections gives the second; this is also the explicit conclusion of
Lemma 59.100.8. This compares coefficients pulled back from the
\'etale site; it does not identify all flat group schemes with
\'etale group schemes. For example, in characteristic $p$ the group
scheme $\alpha_p$ is flat and is not \'etale.

For the second assertion, the divided powers extend along the smooth
lifting. The two pullbacks to the divided power envelope of the
diagonal are canonically identified by the crystal property.
Their first order difference is the connection, and the cocycle on
the triple product makes its curvature zero. The crystalline
Poincar\'e lemma identifies the derived direct image on $Y$ with the
displayed de Rham complex. These are exactly the hypotheses and
conclusion of \cite[Remark 60.24.11, tag 07MW]{BStacks}, whose proof
reduces to Lemma 60.21.4. The converse classification uses
topologically quasi-nilpotent integrable connections on the completed
divided power envelope, not arbitrary integrable connections:
\cite[Proposition 60.17.4, tag 07JH, with conditions 07JB--07JE]{BStacks}.
Thus the stated comparison genuinely includes infinitesimal
coefficients while preserving the needed nilpotence condition.
\end{proof}

\begin{theorem}[Known with checked deduction; subject 13]\label{thm:B3-approx}
Let $(A,I)$ be a henselian pair with $A$ a noetherian G-ring. Given
polynomials $f_1,\ldots,f_m\in A[T_1,\ldots,T_n]$, a solution
$\widehat a\in\widehat A^n$ in the $I$-adic completion, and $N\ge1$,
there is an exact solution $a\in A^n$ with
$a\equiv\widehat a\pmod{I^N\widehat A}$.

If $A=R$ is a henselian discrete valuation ring with uniformizer $t$,
and the fraction field extension $\operatorname{Frac}\widehat R/
\operatorname{Frac}R$ is separable in positive characteristic, there
are constants $c,b\ge1$, depending on the finite polynomial system,
such that every solution modulo $t^{cn+b}$ is congruent modulo $t^n$
to an exact solution in $R$, for every $n\ge1$.
\end{theorem}
\begin{proof}
For the first assertion, the map $A\to\widehat A$ is regular.
Popescu desingularization factors the map
$A[\mathbf T]/(\mathbf f)\to\widehat A$ through a smooth
$A$-algebra $B$. Its reduction modulo $I^N$ lifts to an \'etale
neighbourhood $A'$ of the henselian pair. Henselianity supplies a
retraction $A'\to A$, and the images of the coordinates give $a$.
This is the full argument in
\cite[Lemma 16.14.1, tag 0AH5]{BStacks}, with its cited regularity
and smoothing results. Artin's original theorem has the more
restricted coefficient hypothesis that $A$ is the henselization of
a finite type algebra over a field or an excellent discrete valuation
ring: \cite[Theorems 1.10 and 1.12, p.~26]{BArtin1969}.

The second assertion is \cite[Theorem 1 and Corollary 1,
p.~59]{BGreenberg1966}, with constants enlarged to absorb the
threshold and the additive loss in its exponent. Explicitly, from
the source's precision $\lfloor v/c_0\rfloor-s_0$, choose
$c=c_0$ and $b$ at least both its threshold and $c_0(s_0+1)$.
Taking $v=cn+b$ gives the claimed precision. This explains the
asymptotic statement about Greenberg functors; a linear bound in
arbitrary dimension is not inferred from the first assertion.
\end{proof}

\begin{theorem}[Known and refuted in general; subject 14]\label{thm:B3-topos}
A small site with finite limits whose coverings admit finite
refinements has enough points. An unrestricted Grothendieck topos
need not have even one point.
\end{theorem}
\begin{proof}
The positive assertion is Deligne's theorem, SGA~4 VI, appendix,
Proposition 9.0, with its proof opened as
\cite[Proposition 7.39.3, tag 00YQ]{BStacks}. The original statement
and proof are also opened in
\cite[Expos\'e VI, Proposition 9.0, reedition pp.~464--466]{BSGA4Complete}.
Its conclusion is the
existence of a conservative family of inverse-image functors to sets,
which is exactly the fibre-functor interpretation of points.

Here is a complete counterexample to dropping the finiteness
hypothesis. Let $B$ be the complete Boolean algebra of regular open
subsets of $\R$, regarded as a locale; joins are interiors of closures
of unions and complements are interiors of set complements. It is
atomless because every nonempty regular open contains two disjoint
nonempty open intervals. A point of its sheaf topos would induce a
frame homomorphism $p:B\to\{0,1\}$ preserving arbitrary joins and
finite meets. Put $a=\bigvee\{b:p(b)=0\}$. Then $p(a)=0$, so
$c=\neg a$ is nonzero. Any nonzero $b\le c$ must satisfy $p(b)=1$,
since otherwise $b\le a\wedge c=0$. Atomlessness gives
$0<b<c$. Both $b$ and $c\wedge\neg b$ then map to $1$, although
their meet is zero. This contradicts preservation of meets.
Thus $\operatorname{Sh}(B)$ has no point. It is a nondegenerate topos
since $0\ne1$ in $B$. Representability of a stack is a separate
condition; the existence of points of its underlying topos alone
does not assert representability.
\end{proof}

\begin{theorem}[Known with checked deduction; subject 17]\label{thm:B3-picard}
Let $S$ be the spectrum of a discrete valuation ring with perfect
residue field, and let $f:C\to S$ be a proper flat relative curve.
Assume $C$ is regular and $f_*\mathcal O_C=\mathcal O_S$.
Let $\delta$ be the greatest common divisor of the multiplicities
of the irreducible components of the geometric special fibre.
Then $\operatorname{Pic}^0_{C/S}$ is representable if and only if it
is separated, and these conditions hold if and only if $\delta=1$.
When these conditions hold, it is a smooth separated group scheme.
\end{theorem}
\begin{proof}
Regularity implies condition (N) in
\cite[Definition 6.1.4 and Examples 6.1.6]{BRaynaud1970}: the
special fibre has no embedded components, and the local rings of
$C$ at its generic points are discrete valuation rings.
The assumption on $f_*\mathcal O_C$ gives (N)*.
The local rings after strict henselization remain regular and hence
factorial. Over a perfect residue field Raynaud's numbers $\delta$
and $d'$ coincide. Consequently all eight conditions of
\cite[Theorem 8.2.1, pp.~66--67]{BRaynaud1970} are equivalent;
these include $\delta=1$, representability, and separatedness of
$\operatorname{Pic}^0$. In the representable case the source's
formal smoothness statement for relative curves gives smoothness.
This identifies precisely one substantial scope of the manuscript's
margin ``r\'esolu par Raynaud''. Without the stated regularity and
residue-field conditions Raynaud's theorem is a chain of implications,
not eight unconditional equivalences; his Remark 8.2.2 and Section 9
give the exceptions.
\end{proof}

\begin{theorem}[Known with checked deduction; subject 18]\label{thm:B3-descent}
Affine schemes satisfy effective fpqc descent over every base scheme.
An \'etale equivalence relation $R\rightrightarrows U$ of schemes
has an algebraic-space quotient $U/R$, and $U\to U/R$ is \'etale
surjective. Effective flat descent fails in the category of all
schemes, even for descent data on a projective morphism relative to
an \'etale covering.
\end{theorem}
\begin{proof}
For a faithfully flat affine cover $\operatorname{Spec}A\to
\operatorname{Spec}R$, descent of an affine scheme is descent of an
$A$-algebra $B$. Faithfully flat descent gives the equalizer module
$C\subset B$ and an isomorphism $A\otimes_R C\simeq B$.
Compatibility of the cocycle with multiplication makes $C$ an
$R$-algebra, so $\operatorname{Spec}C$ is the descended scheme.
General fpqc covers reduce to this case:
\cite[Lemma 35.37.1, tag 0245]{BStacks}.

The quotient assertion is
\cite[Theorem 65.10.5, tag 02WW]{BStacks}; its proof constructs the
quotient sheaf and its \'etale presentation. The negative assertion
is the explicit Hironaka construction in
\cite[Lemma 110.68.1, tag 08KE]{BStacks}. It starts from the
involution $[x:y:z:w]\mapsto[y:x:w:z]$ on $\mathbf P^3_\C$, removes
the two fixed lines, and blows up the two exchanged conics in
opposite orders over complementary open sets. The resulting
descent datum is projective over an \'etale cover, but a scheme
quotient would supply a divisor with positive degree on a
rationally zero effective curve, a contradiction. Its quotient
exists as a smooth separated algebraic space. Thus replacing
``scheme'' by ``algebraic space'' in this example is substantive,
and not a change in terminology for an already existing scheme.
\end{proof}

\begin{theorem}[Known and refuted without a fibre functor; subject 22]\label{thm:B3-tannaka}
Let $k$ be a field and $\mathcal T$ a rigid abelian symmetric
$k$-linear tensor category with $\operatorname{End}(\mathbf1)=k$.
If $\omega:\mathcal T\to\operatorname{Vect}_k$ is exact, faithful,
and symmetric tensor, then $G=\operatorname{Aut}^{\otimes}(\omega)$
is an affine group scheme and
$\mathcal T\simeq\operatorname{Rep}_k(G)$ as tensor categories.
Rigidity and the abelian tensor axioms alone do not imply this
conclusion, even in characteristic zero.
\end{theorem}
\begin{proof}
The reconstruction is \cite[Theorem 2.11, pp.~20--28]{BDeligneMilne}.
The coalgebra is reconstructed from the exact faithful functor;
the tensor product supplies multiplication and rigidity the
antipode. These are exactly the hypotheses of the theorem, rather
than an assumed fibre functor constructed from rigidity.
An affine group scheme over a field is an inverse limit of its
finite type group-scheme quotients, which gives the proalgebraic
description in the manuscript.

For the negative assertion take finite-dimensional super vector
spaces over a characteristic-zero field, with the sign symmetry.
This category is rigid, abelian, and has unit endomorphism ring $k$.
The odd line $L$ has categorical dimension $-1$: the symmetry in
the evaluation--coevaluation trace contributes one minus sign.
A symmetric tensor functor preserves that trace, whereas the
dimension of a nonzero ordinary finite-dimensional vector space
is a positive integer. Faithfulness excludes sending $L$ to zero.
Hence no such $\omega$ exists. The sign in subject 22 therefore
cannot be ignored. In the neutral case the stack of fibre functors
is the gerbe of $G$-torsors by
\cite[Theorem 3.2 and Example 3.8]{BDeligneMilne}; its representation
category recovers $\mathcal T$. The precise nonneutral assertion is
Theorem~\ref{thm:B3-nonneutral} below.
\end{proof}

\begin{theorem}[Known with checked deduction; subject 28]\label{thm:B3-motives}
For every finite type $\Z$-scheme $X$ there is a rational tensor
triangulated category of Beilinson motives $\mathrm{DM}_{B}(X)$.
Its constructible subcategory has the six operations, with
$f_!$ attached to separated finite type morphisms, and Verdier
duality for separated finite type $X\to\operatorname{Spec}\Z$.
An explicit categorical construction is
\[
 \mathrm{DM}_{B}(X)\simeq
 \operatorname{Ho}(H_{B,X}\text{-}\operatorname{Mod}),
\]
where $H_{B,X}$ is the Beilinson motivic-cohomology spectrum,
the degree-zero Adams summand of the rational $KGL$ spectrum
representing homotopy invariant algebraic $K$-theory.
\end{theorem}
\begin{proof}
The displayed equivalence is
\cite[Theorem 14.2.9]{BCisinskiDeglise}; Corollary 14.2.11 gives the
motivic formalism. Finite type schemes over $\Z$ are noetherian
excellent of finite dimension. Consequently every condition in
\cite[Theorem 15.2.1]{BCisinskiDeglise} needed for the six operations
preserving constructibility holds in the stated scope. For duality,
take $B=S=\operatorname{Spec}\Z$ in Theorem 15.2.4: $S$ is regular,
excellent, and one-dimensional. If $a:X\to S$ is separated finite
type, then $a^!\mathbf1_S$ is dualizing and
\[
 M\simeq R\mathcal Hom_X(
 R\mathcal Hom_X(M,a^!\mathbf1_S),a^!\mathbf1_S)
 \quad(M\in\mathrm{DM}_{B,c}(X)).
\]
All maps and categories in this construction are defined without
assuming the standard conjectures. It answers the requested
categorical construction in a triangulated, rational form. It does
not construct a conjectural abelian category of mixed motives or
a motivic $t$-structure; that stronger interpretation is not a
consequence of the cited theorem.
\end{proof}

\begin{theorem}[Known with checked deduction; subject 2]\label{thm:B3-flat-duality}
Let $K$ be a complete discretely valued field of characteristic $p>0$
with finite residue field, and $N$ a finite commutative $K$-group
scheme. Write $N^D=\mathcal Hom(N,\mathbf G_m)$ for its Cartier dual.
For every $r$, cup product and the invariant of the Brauer group give
Pontryagin dualities
\[
 H^r_{\mathrm{fppf}}(K,N^D)\times
 H^{2-r}_{\mathrm{fppf}}(K,N)\longrightarrow\Q/\Z.
\]
The natural locally compact topologies must be retained.

If $X$ is a smooth proper geometrically connected curve over a finite
field, $U\subseteq X$ is a nonempty open subscheme, and $N/U$ is a
finite flat commutative group scheme, there are canonical
isomorphisms
\[
 H^{3-r}_{c,\mathrm{fppf}}(U,N)
 \simeq\operatorname{Hom}_{\mathrm{cts}}
 (H^r_{\mathrm{fppf}}(U,N^D),\Q/\Z).
\]
Here the group on the right is the dual of the discrete global
cohomology group. Compact support is the mapping-fibre construction
using the completions at $X\setminus U$.
\end{theorem}
\begin{proof}
These are \cite[III, Theorem 6.10, p.~277; Theorem 8.2,
p.~290]{BMilneADT}, with the hypotheses on the field from the
beginnings of Sections 6 and 8. The local proof first handles the
prime-to-$p$ part by ordinary local duality. For the $p$-part it
uses Cartier duality and resolutions by tori in the \'etale case,
and the two differential sequences of Section 5 for the infinitesimal
case; exact sequences and the five lemma finish the d\'evissage.
For global duality, Lemmas 8.3--8.5 reduce to proper curves and
height-one group schemes, where those differential sequences and
Serre duality apply. Thus the statements include infinitesimal
$p$-primary groups, as the dependency on subject 1 requires.
For example, $H^1(K,\alpha_p)=K/K^p$ follows from
$0\to\alpha_p\to\mathbf G_a\xrightarrow{F}\mathbf G_a\to0$.
For $K=\mathbf F_q((t))$ this group is infinite, so describing the
local theorem as a perfect pairing of finite groups would be false.
The global definition has the localization triangle
\[
 R\Gamma_c(U,N)\to R\Gamma(U,N)\to
 \bigoplus_{v\in X\setminus U}R\Gamma(K_v,N)\longrightarrow,
\]
with the conventions of \cite[III, Section 0.6(b)]{BMilneADT}.
This types the support condition rather than assuming that an
arbitrary flat sheaf has the same compact support theory as its
\'etale restriction.
\end{proof}

\begin{theorem}[Known with checked deduction; subject 9]\label{thm:B3-analytic-duality}
Let $X$ be a complex analytic space, countable at infinity and of
bounded dimension. There is a dualizing complex $K_X$ with coherent
cohomology and a trace $H_c^0(X,K_X)\to\C$ such that, for every
coherent $\mathcal O_X$-module $F$, the pairings
\[
 H^q(X,F)\times\operatorname{Ext}^{-q}_c(F,K_X)\to\C,
 \qquad
 H_c^q(X,F)\times\operatorname{Ext}^{-q}(F,K_X)\to\C
\]
are perfect topological dualities after passing to the associated
Hausdorff quotients. The separatedness criteria have a one-degree
shift:
\[
 H^q(X,F)\text{ is Hausdorff}
 \ \Longleftrightarrow\ 
 \operatorname{Ext}^{1-q}_c(F,K_X)\text{ is Hausdorff},
\]
\[
 H_c^q(X,F)\text{ is Hausdorff}
 \ \Longleftrightarrow\ 
 \operatorname{Ext}^{1-q}(F,K_X)\text{ is Hausdorff}.
\]
If $X$ is smooth of dimension $n$, then
$K_X\simeq\omega_X[n]$. In particular, for compact smooth $X$ and
a holomorphic vector bundle $E$, the first pairing is the finite
dimensional perfect pairing
\[
 H^q(X,E)\times H^{n-q}(X,E^\vee\otimes\omega_X)\to\C.
\]
\end{theorem}
\begin{proof}
The two general assertions are
\cite[Section 3, Theorem 1, p.~81; Section 4, Theorem 2,
pp.~82--83]{BRamisRuget}. Their construction embeds local pieces of
$X$ in Stein manifolds, takes the Hom complex from their structure
sheaves into the ambient shifted canonical complexes, and proves
independence of the embedding. The trace-compatible local complexes
glue. The two dualities follow by applying the local Stein duality
to a locally finite countable Stein covering and transposing its
\v Cech complex. This explains both the support exchange and the
Hausdorff qualification in the exact source statements.
For smooth $X$ the construction gives $\omega_X[n]$.
For locally free $E$, derived Hom becomes
$E^\vee\otimes\omega_X[n]$, so its degree $-q$ hypercohomology is
$H^{n-q}(X,E^\vee\otimes\omega_X)$. Compactness identifies ordinary
and compact supports, and coherent cohomology is finite dimensional;
hence its canonical topology is Hausdorff. The displayed special
case is therefore an actual consequence of the singular analytic
duality theorem, with every shift accounted for.
\end{proof}

\begin{theorem}[Known with checked deduction; subjects 11--12]\label{thm:B3-neron}
Let $R$ be a discrete valuation ring, $K=\operatorname{Frac}R$, and
$K^{\mathrm{sh}}$ the fraction field of a strict henselization.
A smooth $K$-group scheme $G$ of finite type admits a N\'eron model
of finite type over $R$ if and only if $G(K^{\mathrm{sh}})$ is
bounded. In particular every abelian variety over $K$ has such a
model. More generally an abelian variety over the function field
of a connected Dedekind scheme $S$ has a N\'eron model over $S$.
In each assertion the model represents, on smooth $S$-schemes, the
functor
\[
 Y\longmapsto\operatorname{Hom}_K(Y_K,G):\quad
 \operatorname{Hom}_S(Y,\mathcal G)
 \xrightarrow{\sim}\operatorname{Hom}_K(Y_K,G).
\]
It is unique up to a unique isomorphism respecting the generic fibre.
\end{theorem}
\begin{proof}
The local criterion is \cite[Section 1.3, Theorem 1, p.~16]{BBLR}.
Boundedness means that on a finite affine covering the points can be
partitioned into sets on which all chosen affine coordinates have
uniformly bounded valuation below; it is independent of coordinates
by \cite[Section 1.1, Propositions 3--7]{BBLR}.
Proper varieties have bounded sets of points, so Corollary 2 of
Section 1.3 gives the abelian case. Section 1.4, Theorem 3 glues the
local abelian models to the global one: after removing finitely many
closed points the abelian variety extends to an abelian scheme,
which already has the mapping property, and the local models agree
there by that property. The stated bijection is the N\'eron mapping
property, Definition 1 of Section 1.2. Applying it in both directions
to two models gives inverse maps, since a map extending the identity
is unique. Applying it to multiplication, identity and inverse on
the generic fibre gives the group structure, as in Proposition 6
of that section.

This also identifies the scope of the functor mentioned under
subject 11. Finite type cannot be omitted from the boundedness
criterion: $\mathbf G_{m,K}$ has points with arbitrarily large
positive and negative valuations and therefore has no finite type
N\'eron model, although it has a N\'eron model locally of finite type.
The manuscript's additional phrase ``sch\'emas morcel\'es'' supplies
no definition on these four sheets. The theorem settles the named
N\'eron functor and does not identify every possible meaning of that
undeveloped category.
\end{proof}

\begin{theorem}[Known with checked deduction; nonneutral subject 22]\label{thm:B3-nonneutral}
For the tensor category $\mathcal T$ of
Theorem~\ref{thm:B3-tannaka}, suppose there is an exact symmetric
$k$-linear tensor functor $\omega$ to vector bundles on a nonempty
$k$-scheme $S$. Then
\[
 G=\operatorname{Isom}^{\otimes}_k
       (\operatorname{pr}_2^*\omega,\operatorname{pr}_1^*\omega)
 \longrightarrow S\times_k S
\]
is an affine faithfully flat groupoid and
$\mathcal T\simeq\operatorname{Rep}(S:G)$.
The fpqc stack of fibre functors is a gerbe, and its vector-bundle
representations recover $\mathcal T$.
In characteristic zero such a fibre functor exists over some
nonempty $S$ if and only if every object has categorical dimension
a nonnegative integer; equivalently, for every object $M$ some
exterior power $\bigwedge^nM$ vanishes.
\end{theorem}
\begin{proof}
The first assertion, including faithful flatness rather than mere
representability of the Isom functor, is exactly
\cite[Theorem 1.12, p.~116]{BDeligne1990}. The categorical hypotheses
are those in its Definition 1.2; its Definition 1.9 makes the target
locally free of finite rank automatically. Given two fibre functors,
apply the theorem to their disjoint union over $S_1\amalg S_2$.
The off-diagonal restriction is the faithfully flat Isom scheme,
so the functors become isomorphic fpqc locally. This is the argument
of Remark 1.13(a), and proves the gerbe assertion. Descent for vector
bundles on this presentation identifies its representations with
$\operatorname{Rep}(S:G)$, hence with $\mathcal T$.
The intrinsic characteristic-zero criterion is
\cite[Theorem 7.1, p.~165]{BDeligne1990}. Its hypotheses exclude the
odd super line in Theorem~\ref{thm:B3-tannaka}, whose categorical
dimension is $-1$ and whose exterior powers do not eventually vanish.
Deligne explicitly repairs the incomplete fibre-functor comparison
argument in Saavedra's treatment (Introduction and Section 1.10);
this is why a reference to Saavedra alone without those hypotheses
would not establish the general reconstruction asserted here.
\end{proof}

\begin{theorem}[Known with checked deduction; subject 26]\label{thm:B3-rosenlicht}
Let $S$ be any scheme and $\overline C\to S$ a flat finitely
presented proper relative curve with geometrically integral fibres,
locally projective over $S$. Let $T\subset\overline C$ be a relative
effective Cartier divisor, flat over $S$, such that
$C=\overline C\setminus T$ is smooth over $S$.
There is an explicitly constructed inverse limit of smooth
commutative $S$-group schemes $J^\bullet_\infty$, with a morphism
$\varphi:C\to J^\bullet_\infty$, such that, for every smooth
commutative $S$-group scheme $G$, composition gives a bijection
\[
 \operatorname{Hom}_{S\text{-}\mathrm{gp}}(J^\bullet_\infty,G)
 \xrightarrow{\sim}\operatorname{Hom}_S(C,G).
\]
Thus no chosen section of $C/S$ is necessary; the degree components
are part of the representing group.
\end{theorem}
\begin{proof}
This is \cite[Theorem 1.4, p.~204, with hypotheses in Section 0,
p.~203]{BContou1979}. For clarity its construction is as follows.
Put $T_n=V(\mathcal I_T^{n+1})$ and pinch $T_n$ to $S$ in
$\overline C$, producing $\overline C_{(n)}$ with its distinguished
section. The corresponding rigidified Picard group $J^\bullet_n$
has a degree map to $\Z_S$ and exact sequence
\[
 0\longrightarrow J^0_n\longrightarrow J^\bullet_n
 \longrightarrow\Z_S\longrightarrow0.
\]
The transition maps are affine, so their inverse limit exists as an
$S$-scheme. A relative point of $C$ gives a Cartier divisor away
from the pinched section, hence compatible degree-one Picard classes
and the required $\varphi$. These are the constructions of Section 1,
including formulas (1.1)--(1.3).
The proof reduces by finite presentation to an affine noetherian
base. Over a reduced base a sufficiently high symmetric power of
$C$ covers a degree component faithfully flatly; the field-level
Rosenlicht factorization descends along it. For a nilpotent
thickening, the obstruction to lifting a group homomorphism is the
extension by the vector-group functor associated with the square-zero
ideal. Proposition 3.2 and Lemma 3.5 kill it after increasing $n$.
Induction on the nilradical gives existence and uniqueness over the
whole base. This accounts for the nonreduced and arbitrary-base
clauses rather than reducing the assertion only to geometric fibres.
The inverse limit and all degree components are essential: a single
fixed generalized Jacobian with a fixed modulus need not factor
morphisms of arbitrarily large modulus.
\end{proof}

\begin{theorem}[Known bound and explicit obstruction; subject 3]\label{thm:B3-pprimary}
Let $Y$ be a smooth projective connected curve of genus $g$ over an
algebraically closed field $k$ of characteristic $p>0$, and let $N$
be a constructible $\mathbf F_p$-sheaf whose sections have open support.
There are explicitly defined nonnegative rational local indices $C_y(N)$,
zero except at finitely many points, such that
\[
 \chi(Y,N)\ge (1-g)\operatorname{rank}N-\sum_{y\in Y}C_y(N).
\]
An equality determined solely by the genus, rank and multiset of local
ramification representations is false for $p$-coefficients.
Already over $k=\overline{\mathbf F}_7$, the two degree-two covers
\[
 E_0:y^2=x^3-x,\qquad E_1:y^2=x^3-x+1,
 \qquad\beta_i:E_i\longrightarrow\mathbf P^1,
\]
give rank-two sheaves $N_i=\beta_{i*}\mathbf F_7$ with four identical
tame local representation types and
$\chi(\mathbf P^1,N_0)=1$, $\chi(\mathbf P^1,N_1)=0$.
The branch points themselves occupy different positions; the assertion
concerns local types and their multiplicities, not an isomorphism of
sheaves with a fixed labelled branch divisor.
\end{theorem}
\begin{proof}
The bound is \cite[Theorem 1.0.7 and Corollary 3.1.18]{BMiller2010}.
Here is the explicit invariant and deduction. Put
$M=\mathcal Hom_{\mathbf F_p}(N,\mathcal O_Y)$ with its Frobenius
operator. Its minimal root $M_0$ is the smallest coherent submodule
whose Frobenius iterates generate $M$ and for which
$M_0\subset M_1=\mathcal O_Y F(M_0)$.
On a smooth curve it is locally free, $M_1\simeq F_Y^*M_0$, and
\[
 C_y(N)=\frac{\operatorname{length}_{\mathcal O_{Y,y}}
                (M_1/M_0)_y}{p-1}.
\]
Existence and uniqueness are proved in Section 2.2. The
Riemann--Hilbert construction gives the exact sequence
$0\to N\to M_0^\vee\xrightarrow{1-F}M_0^\vee\to0$
(Section 3.1). Its cohomology compares Frobenius-fixed subspaces with
coherent cohomology and yields
$\chi(Y,N)\ge\chi(Y,M_0^\vee)$.
Since $\deg(M_1)=p\deg(M_0)$, the sum of the displayed lengths is
$(p-1)\deg(M_0)$. Riemann--Roch therefore gives exactly the bound.

We verify the counterexample without an unspecified generic elliptic
curve. For $y^2=x^3+ax+b$ in characteristic seven, the smoothness
condition is $4a^3+27b^2\ne0$. It holds for $(a,b)=(-1,0)$ and
$(-1,1)$, giving respectively $3$ and $2$ modulo seven. Both $x$-maps
are finite separable double covers ramified simply at the three roots
of the cubic and at infinity. Each strict local ramification
representation is the rank-two permutation representation of a group
of order two, so the four local types agree up to identifying their
strict local parameters. Away from the branch points each sheaf is
locally constant of rank two.

For $\eta=dx/y$, the Cartier operator is
$C(\eta)=c_6^{1/7}\eta$, where $c_6$ is the coefficient of $x^6$ in
$(x^3+ax+b)^3$. Indeed
$\eta=y^{-7}(x^3+ax+b)^3dx$, and Cartier retains precisely exponents
congruent to $6$ modulo seven. That coefficient is $3b$.
Thus Frobenius on $H^1(E_0,\mathcal O_{E_0})$, dual to Cartier,
is zero, whereas on the one-dimensional space for $E_1$ it is
bijective and semilinear. The Artin--Schreier sequence gives
\[
 H^1(E_i,\mathbf F_7)=
 \ker(F-1:H^1(E_i,\mathcal O_{E_i})\to H^1(E_i,\mathcal O_{E_i})).
\]
The dimensions are respectively zero and one. The cokernel of $F-1$
is zero in both cases: for zero $F$ the map is $-1$, and otherwise
it is a nonconstant additive polynomial on the algebraically closed
field. Hence $H^2(E_i,\mathbf F_7)=0$, and both $H^0$ have dimension
one. Finite proper base change gives
$R\beta_{i*}\mathbf F_7=\beta_{i*}\mathbf F_7$, since its geometric
fibres are finite schemes with no positive-degree cohomology.
This proves the asserted Euler numbers. It is a completely specified
version of \cite[Example 3.2.9]{BMiller2010}. The polynomial
coefficients and discriminants are checked in the accompanying code.
For these examples each $C_y=1/2$, so the bound is $\chi\ge0$;
it is sharp for $E_1$ and strict for $E_0$.
\end{proof}

\begin{theorem}[Known with checked deduction; subjects 4 and 6]\label{thm:B3-trace}
Let $X$ be a proper scheme of finite type over an algebraically closed
field $k$, let $\ell\ne\operatorname{char}k$, and let
$c=(c_1,c_2):C\to X\times_kX$ be a correspondence with $C$ proper.
For $F\in D^b_c(X,\Q_\ell)$ and a cohomological correspondence
$u:c_{2!}c_1^*F\to F$, the Lefschetz--Verdier formula is
\[
 \operatorname{Tr}(R\Gamma(u))
   =\sum_{\beta\in\pi_0(\operatorname{Fix}(c))}\operatorname{LT}_\beta(u).
\]
The local terms are canonically defined by the trace class in the
dualizing complex of the fixed-point scheme.
Consequently a finite \'etale map $f:Y\to X$ of constant degree $n$
satisfies
\[
 \chi(Y,\Q_\ell)=n\chi(X,\Q_\ell).
\]
No smoothness assumption on $X$ is needed. Properness is essential.
\end{theorem}
\begin{proof}
The trace class is the composite, in the notation
$\Delta':\operatorname{Fix}(c)\to C$,
\[
 R\mathcal Hom(c_1^*F,c_2^!F)
 \longrightarrow c^!(D_XF\boxtimes F)
 \longrightarrow c^!\Delta_*K_X
 \simeq\Delta'_*K_{\operatorname{Fix}(c)}.
\]
The first arrow is the duality identification, the second is evaluation,
and the third is proper base change for the diagonal. Restricting the
induced degree-zero trace class to a proper connected component and
integrating defines $\operatorname{LT}_\beta$.
Compatibility with proper pushforward is
\cite[Proposition 1.2.5, proved in Section 4.3]{BVarshavsky2005}.
Push the correspondence to a point: evaluation then becomes the
ordinary alternating trace on cohomology. This is the proof of
\cite[Corollary 1.2.6]{BVarshavsky2005}, so it gives the displayed
formula with its actual local terms, not an unjustified replacement
by traces on stalks.
For finite coefficients, the original proper compatibility theorem
and its trace specializations are also opened in
\cite[Expos\'e III, Theorem 4.4, p.~102;
Corollaries 4.7--4.8, p.~109]{BSGA5}.

For the Euler formula work one connected component and one connected
finite cover at a time. Dominate the cover by a finite \'etale Galois
cover $Z\to X$ with group $G$, and write $Y=Z/H$. Every nonidentity
element of $G$ acts on $Z$ without fixed points. The trace formula
therefore gives $\operatorname{Tr}(g\mid R\Gamma(Z,\Q_\ell))=0$
for $g\ne1$. Invariants under a finite group are exact over
$\Q_\ell$, and finite descent gives
$R\Gamma(Z/G,\Q_\ell)=R\Gamma(Z,\Q_\ell)^G$.
Averaging traces yields
$\chi(X)=\chi(Z)/|G|$ and $\chi(Y)=\chi(Z)/|H|$.
Their ratio is $[G:H]=n$. Additivity over components proves the general
case. In characteristic $p$, the map
$\mathbf A^1\to\mathbf A^1$, $t\mapsto t^p-t$, is finite \'etale of
degree $p$ because its derivative is $-1$. Both compactly supported
Euler characteristics are one, so the corresponding assertion for
arbitrary nonproper schemes would be false.
\end{proof}

\begin{theorem}[Known with checked deduction; higher-dimensional part of subject 4]\label{thm:B3-index}
Let $X$ be a smooth projective variety of dimension $d$ over an
algebraically closed field, and let $F$ be a constructible sheaf of
vector spaces over a finite field $\Lambda$ of characteristic
invertible in the base field. Its characteristic cycle $CC(F)$ is an
integral conical $d$-cycle on $T^*X$, characterized by the Milnor
formula for vanishing cycles, and
\[
 \chi(X,F)=(CC(F),T_X^*X)_{T^*X}.
\]
For a locally constant sheaf of rank $r$,
$CC(F)=(-1)^d r[T_X^*X]$. On a curve the additional fibre
coefficients are the negatives of the Artin conductors.
\end{theorem}
\begin{proof}
Existence and the Milnor characterization are
\cite[Theorem 4.9]{BSaito}; the constant-sheaf normalization and the
curve computation are Lemma 4.11(1),(3), with the curve formula
in equation (4.19). The index formula is exactly
Theorem 6.13. All its hypotheses hold in the stated scope. Thus the
formula supplies a precise higher-dimensional replacement for an
Euler formula in terms of local ramification data: those data assemble
into a cycle in the cotangent bundle and are paired with the zero
section. It does not assert that they are a single ordinary Artin
representation of a single group, a notion the manuscript does not
specify in higher dimension.
\end{proof}

\begin{theorem}[Known with checked deduction; subject 5]\label{thm:B3-operators}
Let $C$ be a smooth projective connected curve over an algebraically
closed field of characteristic $p\ge0$, let $R$ be an arbitrary ring,
and let $F$ be a constructible \'etale sheaf of finite $R$-modules
killed by a power of a prime $\ell\ne p$. In the Grothendieck group
$K_0(\mathrm{FinMod}_R)$,
\[
 \sum_i(-1)^i[H^i(C,F)]
   =(2-2g)[F_{\bar\eta}]-\sum_x\epsilon_x^R(F),
\]
where
\[
 \epsilon_x^R(F)=
 [F_{\bar\eta}]-[F_{\bar x}]
 +[\operatorname{Hom}_{\Z_\ell[G_x]}(P_{G_x},F_{\bar\eta})].
\]
Here $G_x$ is a finite local monodromy quotient and $P_{G_x}$ its
projective Swan module. This definition is independent of the
chosen quotient. It applies in particular to the constructible
extension of the $\ell^m$-torsion of an abelian variety over $k(C)$,
with any ring of endomorphisms acting on it.
\end{theorem}
\begin{proof}
This is the operator-valued theorem of Grothendieck proved in
\cite[Part I, Theorem 1, p.~133]{BRaynaud1966}. Part I, Section 1
constructs the Swan module and proves independence under enlarging
the splitting extension. Projectivity makes its Hom functor exact;
therefore the conductor class is additive in exact sequences.
Section 2, formulas (1)--(1 ter), gives the displayed expression,
first for $F=i_*F_\eta$ and then for every constructible $F$ by the
skyscraper kernel and cokernel of $F\to i_*F_\eta$.
The proof in Section 4 uses induction from finite covering curves,
Riemann--Hurwitz, and additivity, without forgetting the $R$-action.
For an abelian variety $A/k(C)$ choose the open of good reduction
$j:U\hookrightarrow C$. Its abelian scheme $\mathcal A/U$ has finite
\'etale torsion $\mathcal A[\ell^m]$, so
$F=j_*\mathcal A[\ell^m]$ satisfies every hypothesis.
An endomorphism of $A$ extends on its N\'eron model and preserves this
sheaf, giving the specified $R$-action. The Kummer sequences
$0\to\mathcal A[\ell^m]\to\mathcal A\xrightarrow{\ell^m}
\mathcal A\to0$ on $U$ identify these finite-level terms with the
corresponding derived torsion of abelian cohomology. This is the
operator-preserving content; replacing it by dimensions alone would
lose the information requested in subject 5. The restriction
$\ell\ne p$ is essential, and subject 3 is treated separately.
\end{proof}

\begin{theorem}[Known with checked deduction; subject 7]\label{thm:B3-arithmetic-duality}
Let $B$ be the spectrum of a ring of integers in a number field,
$f:X\to B$ separated of finite type, and $F$ a torsion \'etale sheaf
on $X$. Bloch's cycle complex $\Z_X^c$ gives a canonical duality
\[
 R\operatorname{Hom}_X(F,\Z_X^c)
  \simeq R\operatorname{Hom}_{\Z}
       (R\Gamma_c(X_{\mathrm{\acute et}},F),\Z)[-1].
\]
For constructible $F$ it induces perfect pairings of finite groups
\[
 \operatorname{Ext}^{2-i}_X(F,\Z_X^c)
     \times H_c^i(X_{\mathrm{\acute et}},F)\longrightarrow\Q/\Z.
\]
Compact support incorporates the modified cohomology at real places;
this modification matters for $2$-primary coefficients.
\end{theorem}
\begin{proof}
The assertions are \cite[Theorem 7.8 and Section 7.2]{BGeisser}.
The source's global convention (p.~4) makes every scheme over a base
separated and of finite type, as explicitly imposed here.
Its Section 7 assumes Beilinson--Lichtenbaum. That input is now a
theorem: the exact comparison needed there is
\cite[Theorem 6.17]{BVoevodsky2011}, together with the
positive-characteristic comparison already used in Geisser's
Sections 2--3. Consequently Theorem 7.3 identifies
$R\operatorname{Hom}_X(F,\Z_X^c)$ with
$R\operatorname{Hom}_B(Rf_!F,\mathbf G_m)[1]$.
Artin--Verdier duality on $B$ identifies this with
$R\operatorname{Hom}_{\Z}(R\Gamma_c(X,F),\Q/\Z)[-2]$.
For a torsion complex the exact sequence
$0\to\Z\to\Q\to\Q/\Z\to0$ identifies the latter with the
$\Z$-dual shifted by $-1$, proving the claimed formula and degree
$2-i$ pairing. Constructibility supplies the finiteness needed for
perfect finite-group duality. This gives a cycle-complex extension of
Artin--Verdier to arbitrary-dimensional finite type arithmetic schemes,
including singular $X$; it is not restricted to regular arithmetic
surfaces.
\end{proof}

\begin{theorem}[Known with checked deduction; subject 19]\label{thm:B3-positive-cones}
Let $X$ be a projective variety over an algebraically closed field.
For every $k$ the pliant cone $\mathrm{PL}^k(X)\subset N^k(X)_\R$,
defined as the closed cone generated by products of pullbacks of
Schubert classes from Grassmannians, is full dimensional and contains
no line. Complete intersections of ample divisors lie in its interior.
It is contained in the nef cone, is preserved by pullbacks, and its
classes and their products cap $[X]$ to pseudoeffective dimensional
cycle classes. Here $N^k(X)_\R$ is the dual of the numerical group
$N_k(X)_\R$; on singular $X$ the cap map to $N_{\dim X-k}(X)_\R$
need not be an isomorphism. On smooth $X$ there are functorial inclusions
\[
 \mathrm{PL}^k(X)\subseteq\mathrm{BPF}^k(X)
 \subseteq\mathrm{Upsef}^k(X)\subseteq\mathrm{Nef}^k(X),
\]
where $\mathrm{BPF}$ is generated by flat moving families and
$\mathrm{Upsef}$ consists of classes whose pullback to every projective
variety is pseudoeffective. Also the cone of pseudoeffective
$k$-dimensional cycles is full dimensional and salient; degree against
an ample $h^k$ is the restriction of a norm on $N_k(X)_\R$.
\end{theorem}
\begin{proof}
The first set of properties is
\cite[Theorem 1.3, proved by Lemmas 3.5, 3.7, 3.14]{BFulgerLehmann}.
The cone construction is Definition 3.2, equivalently the
Grassmannian description in Section 1.1. The inclusions and
basepoint-free and universal pseudoeffectivity properties are
Theorems 1.7 and 1.8. The assertions about dimensional cycles and
ample degree are Theorem 1.4. The degree statement applies to a
nonzero limiting pseudoeffective class as well as to an effective
cycle, so it is stronger than merely observing that a subvariety has
positive degree. These constructions give precise higher-codimension
positivity theories with the formal properties sought in subject 19.
They are not a claim that the several divisor characterizations of
ampleness remain equivalent in higher codimension; the source gives
strict inclusions (Example 1.10) and explains their geometric causes.
\end{proof}

\begin{theorem}[Known with checked deduction; abelian $p$-primary part of subjects 2--3]\label{thm:B3-abelian-duality}
Let $K$ be a complete discretely valued field with perfect residue field
$k$ of characteristic $p>0$, and let $A/K$ be an abelian variety with
dual $A^\vee$. Regard flat cohomology as a complex of sheaves on the
ind-rational pro-\'etale site of $k$, and write
$C^{\mathrm{SD}}=R\mathcal Hom(C,\Z)$. Then the Poincar\'e
biextension induces an isomorphism
\[
 R\underline{\Gamma}(K,A^\vee)^{\mathrm{SDSD}}
 \simeq R\underline{\Gamma}(K,A)^{\mathrm{SD}}[1].
\]
In particular Grothendieck's pairing on the two N\'eron component
groups is perfect, including its $p$-primary part.

For the global version let $X/k$ be a smooth proper geometrically
connected curve over a perfect field of characteristic $p>0$, let
$A/k(X)$ be an abelian variety, and let $\mathcal A$ and
$\mathcal A^\vee_0$ denote respectively its N\'eron model and the
identity component of the dual model. Put
\[
 V=T\mathcal H^1(X,\mathcal A^\vee_0)_{\rm div}\otimes\Q.
\]
Here the subscript means the maximal divisible subsheaf, and $T$ is
its full Tate module. There is a canonical distinguished triangle
\[
 V[-1]\longrightarrow
 R\underline{\Gamma}(X,\mathcal A^\vee_0)^{\mathrm{SDSD}}
 \longrightarrow R\underline{\Gamma}(X,\mathcal A)^{\mathrm{SD}}
 \longrightarrow V.
\]
The term $V$ cannot be discarded in the theorem over arbitrary perfect
$k$.
\end{theorem}
\begin{proof}
The local assertion is \cite[Theorems A and C, pp.~3--4]{BSuzuki2014}.
The field may have equal or mixed characteristic; perfection of the
residue field is a real hypothesis. Section 1.1 defines precisely the
site: its objects are filtered unions of finite products of perfections
of finitely generated fields over $k$. Bold cohomology means the
sheaf-valued functor there, not the discrete cohomology of the single
field $K$. Theorem B identifies the displayed isomorphism with
Grothendieck component duality together with Shafarevich duality, and
proves its descent along finite Galois extensions. Semistable reduction
and the established semistable duality then give the assertion for
all $A$. Thus an older statement of the $p$-primary component pairing
as conjectural is no longer its status in this scope.

The global triangle is \cite[Theorem A, Section 1.2, p.~3]{BSuzuki2020}.
There $R\underline{\Gamma}(X,-)$ is the derived pro-\'etale
sheafification of $k'\mapsto\Gamma(X\times_k k',-)$. The
biextension and the degree map of the Picard scheme give a pairing
with target $\Z$ and hence the middle arrow. Sections 2--3 combine
finite flat global duality, the just-stated local theorem, the height
pairing and crystalline finiteness to identify its cone with the
explicit $V$. Theorem 3.4.1 gives its cohomological consequences.
The paper explains immediately after Theorem A that $V$ can be
nonzero. Consequently neither a perfect pairing of finite discrete
groups in all degrees nor a naive numerical Euler formula follows
by simply deleting it. These statements retain the group-scheme and
$p$-primary information of the original programme.
\end{proof}

\begin{theorem}[Known specialization and explicit failure of full invariance; subject 16]\label{thm:B3-specialization}
For a proper smooth morphism with geometrically connected fibres over
a locally noetherian base, specialization between geometric fibres
induces a surjection of \'etale fundamental groups. If the special
residue characteristic is $p>0$, it induces an isomorphism on the
maximal prime-to-$p$ quotients; in residue characteristic zero it is
an isomorphism on the full groups.
The assertion of full invariance in positive characteristic is false,
even for smooth proper elliptic curves.
\end{theorem}
\begin{proof}
The positive assertions are \cite[Expos\'e X, Theorem 3.8 and
Corollary 3.9, original p.~283]{BSGA1}, not a theorem about all
pro-$p$ quotients. To give a complete counterexample put
$k=\overline{\mathbf F}_3$, $R=k[[t]]$ and consider the elliptic scheme
\[
 \mathcal E/R:\quad y^2=x(x-1)(x-\lambda),\qquad\lambda=2+t,
\]
with the usual point at infinity. Its discriminant
$16\lambda^2(1-\lambda)^2$ is a unit in $R$, so the projective
Weierstrass model is smooth and proper with geometrically connected
fibres. For a Legendre equation in characteristic three, Cartier on
$dx/y$ is the cube root of the coefficient of $x^2$ in
$x(x-1)(x-\lambda)$, namely $-(1+\lambda)=-t$.
It is zero at the special fibre and nonzero at the geometric generic
fibre. The Artin--Schreier argument in
Theorem~\ref{thm:B3-pprimary} now gives
\[
 \dim_{\mathbf F_3}H^1(\mathcal E_{\bar s},\mathbf F_3)=0,
 \qquad
 \dim_{\mathbf F_3}H^1(\mathcal E_{\bar\eta},\mathbf F_3)=1.
\]
These groups equal continuous homomorphisms from the respective
fundamental groups to $\Z/3$, so the full groups cannot be isomorphic.
The coefficient identity is checked by the attached exact computation.
This is our specified counterexample, not an unsupported identification
with the unnamed example attributed to Artin in the margin. The
manuscript's undeveloped word ``vrai'' does not define another group;
we state explicitly that the specialization theorem and counterexample
here concern $\pi_1^{\mathrm{\acute et}}$.
\end{proof}

\begin{theorem}[Known with checked deduction; subject 20]\label{thm:B3-homogeneous}
Let $G$ be a connected simply connected semisimple group over an
algebraically closed field, $P\subseteq G$ a parabolic subgroup,
and $X_w\subseteq G/P$ a Schubert variety. For a dominant character
$\lambda$ of $P$, with the convention that $\mathcal L(\lambda)$
is globally generated, one has
\[
 H^i(X_w,\mathcal L(\lambda)|_{X_w})=0\quad(i>0),
\]
and restriction $H^0(G/P,\mathcal L(\lambda))\to
H^0(X_w,\mathcal L(\lambda)|_{X_w})$ is surjective.
In particular every nef line bundle on the complete homogeneous
space $G/P$ has vanishing higher cohomology. This holds in every
characteristic.
\end{theorem}
\begin{proof}
In positive characteristic these are
\cite[Theorem 3.1.1(a)--(b), pp.~85--86]{BBrionKumar}, with the
conventions of Section 2.1. Its proof uses a splitting of $G/P$
compatible with all Schubert varieties and with an ample boundary
(Theorem 2.2.5), followed by the cohomology vanishing theorem for that
splitting (Theorem 1.4.8). This gives the vanishing for semiample
bundles, not only sufficiently high powers. Dominant characters are
exactly the globally generated, equivalently nef, line bundles in
these conventions. In characteristic zero use the split Chevalley
model, the Schubert closure and its line bundle over $\Z$.
After restricting to a dense open of $\operatorname{Spec}\Z$ the
model is flat. Proper semicontinuity makes the locus of nonzero
higher cohomology closed; it misses every positive characteristic
point of that open by the theorem just used, hence also its generic
point. The exact restriction sequence gives surjectivity by the same
argument. Extension from $\overline\Q$ to any algebraically closed
characteristic-zero field follows by flat cohomology base change.
This proves the asserted all-characteristic formulation. The adjective
``complete'' is essential to the meaning of homogeneous space used
here; no vanishing for arbitrary line bundles on arbitrary homogeneous
spaces is asserted.
\end{proof}

\begin{theorem}[Known with explicit constructions; subject 27]\label{thm:B3-adeles}
Let $X$ be noetherian and $F$ a quasi-coherent $\mathcal O_X$-module.
Parshin--Beilinson ad\`eles give a functorial flasque resolution
\[
 0\longrightarrow F\longrightarrow\mathbb A^0(F)
 \longrightarrow\mathbb A^1(F)\longrightarrow\cdots.
\]
Consequently the alternating flag complex of global ad\`eles computes
$H^*(X,F)$. If $X$ is also reduced, excellent and pure dimensional of
dimension $n$, the ad\`ele ring of a saturated length-$n$ flag is a
finite product of $n$-local fields with last residue fields finite over
the residue field of the final point.

For the local-Jacobian part, let $A$ be noetherian,
$S=\operatorname{Spec}A$, $\mathfrak C=\operatorname{Spf}A[[T]]$,
and $U=\operatorname{Spec}A((T))$. The commutative group functor
\[
 \mathcal J(R)=R((T))^\times
\]
on $A$-algebras has a canonical natural bijection
\[
 \operatorname{Hom}_{S\text{-}\mathrm{gp}}(\mathcal J,G)
 \simeq G(U)
\]
for every smooth separated commutative $S$-group scheme $G$.
It is implemented by an Abel--Jacobi morphism
$U\to\mathcal J_{\rm omb}$ and the canonical extensions of group
homomorphisms to this intermediate functor. It is not implemented by
a morphism $U\to\mathcal J$.
\end{theorem}
\begin{proof}
Here is an explicit definition of the adelic object, so the assertion
is more than a name. For coherent $F$, let $K_r$ be a set of strict
specialization flags of length $r+1$, and write $\eta K_r$ for the
remaining flags after fixing the first point $\eta$. Recursively put
\[
 A(K_0,F)=\prod_{\eta\in K_0}\varprojlim_m
       F\otimes_{\mathcal O_X}\mathcal O_{X,\eta}/\mathfrak m_\eta^m,
 \quad
 A(K_r,F)=\prod_{\eta\in X}\varprojlim_m
       A(\eta K_r,F\otimes\mathcal O_{X,\eta}/\mathfrak m_\eta^m).
\]
For quasi-coherent $F$ take the filtered colimit over coherent
submodules. Sheafify these constructions on opens and take the
alternating sum of face maps. The precise recursion, exactness, and
flasque resolution are \cite[Section 2.1, formulas (2.1)--(2.3),
Theorem 2.1]{BBraunlingEtAl}. The theorem there is Beilinson's,
with the complete proof credited to Huber; the statement is opened
in this full primary exposition. Flasque acyclicity proves the
cohomology assertion. The assertion about fields is its Theorem 0.2(1),
proved in Section 4. In dimension one this is the familiar complex
$K(X)\oplus\prod_x\widehat{\mathcal O}_{X,x}\to
\prod'_x\widehat K_x$ for a regular integral curve; the recursive
restrictions specify the higher-dimensional restricted product.

The local-Jacobian theorem is
\cite[Theorem 1.6.6, p.~54]{BContou2013}; the loop functor description
is given in its abstract and Section 0. For the subgroup
$J^0(R)=R[[T]]^\times$, the intermediate universal property is
\[
 \operatorname{Hom}_{S\text{-}\mathrm{gp}}(J^0,G)
 \simeq G(U)/G(\operatorname{Spec}A[[T]])
\]
(Theorem 1.4.4, p.~46). The construction adds the valuation components
and the formal negative Witt part, giving the full loop functor.
Sections 2--5 prove the universal property over fields, then over
complete discrete valuation rings, then reduced noetherian rings,
and finally through nilpotent thickenings. Section 1.5 constructs
$\mathcal J_{\rm omb}$ and its Abel--Jacobi map; Section 1.6 proves
that its composition with the extended homomorphism gives precisely
the displayed bijection. The distinction is necessary even for a
field base: the formal negative Witt factor has only its identity
section on the reduced scheme $U$, so a literal $U\to\mathcal J$
could not detect its nonzero homomorphisms. This is the explicit
obstruction on p.~5 of the source. The theorem is therefore stated
in its actual functor category and with its noetherian hypothesis,
rather than with an incorrectly typed ordinary universal morphism.
\end{proof}

\begin{theorem}[Known with checked deduction; differential-operator part of subject 8]\label{thm:B3-dmodule}
Let $X$ be a complex analytic manifold of dimension $d$, countable
at infinity, and $\mathcal D_X$ its sheaf of holomorphic differential
operators. For every bounded complex $M$ of left $\mathcal D_X$-modules
with holonomic cohomology, put
\[
 \mathrm{DR}(M)=R\mathcal Hom_{\mathcal D_X}(\mathcal O_X,M),
 \qquad
 \mathrm{Sol}(M)=R\mathcal Hom_{\mathcal D_X}(M,\mathcal O_X).
\]
There is a canonical isomorphism
\[
 \mathrm{DR}(M)\simeq
 R\mathcal Hom_{\C_X}(\mathrm{Sol}(M),\C_X).
\]
Equivalently, with Verdier duality
$D_X(-)=R\mathcal Hom_{\C_X}(-,\C_X[2d])$, one has
\[
 \mathrm{DR}(M)[d]\simeq D_X(\mathrm{Sol}(M)[d]).
\]
No regular-singularity assumption is required for this local duality.
\end{theorem}
\begin{proof}
This is the local duality theorem of Mebkhout and Kashiwara--Kawai
in the exact unshifted conventions of
\cite[Theorem 3.1.1, p.~67]{BNarvaez2004}.
Section 2.1 constructs the Spencer resolution
$\mathcal D_X\otimes\bigwedge^\bullet T_X\to\mathcal O_X$;
Lemma 2.1.1 identifies its Hom complex with the ordinary de Rham
complex of an integrable connection, in degrees $0,\ldots,d$.
Section 2.2 defines the displayed morphism by evaluation and the
holomorphic Poincar\'e lemma. Sections 3.3--3.4 verify its
compatibility with both Serre and local analytic duality and prove
it is an isomorphism, including the comparison diagram needed for
the Kashiwara--Kawai proof. Applying a shift of $d$ and substituting
the displayed definition of $D_X$ gives the second formula directly:
$D_X(S[d])=R\mathcal Hom(S,\C_X)[d]$.
Finite rank integrable connections are holonomic, so they satisfy
these hypotheses. Crystals require the separate positive-characteristic
PD and quasi-nilpotence conditions already stated in
Theorem~\ref{thm:B3-flat-crys}; local analytic duality alone is not a
crystalline comparison theorem.
\end{proof}

\begin{theorem}[Known with checked deduction; subject 10]\label{thm:B3-riemann-roch}
For every separated finite type scheme $X$ over a field there is a
Riemann--Roch transformation
\[
 \tau_X:G_0(X)\longrightarrow\mathrm{CH}_*(X)_\Q
\]
commuting with proper pushforward and satisfying
$\tau_X(E\otimes F)=\operatorname{ch}(E)\cap\tau_X(F)$ for vector
bundles $E$. If $X$ is smooth, then
$\tau_X(F)=\operatorname{ch}(F)\operatorname{td}(T_X)\cap[X]$.
Thus for a proper morphism $f:X\to Y$ of smooth schemes,
\[
 f_*\bigl(\operatorname{ch}(E)\operatorname{td}(T_X)\cap[X]\bigr)
 =\operatorname{ch}(Rf_*E)\operatorname{td}(T_Y)\cap[Y].
\]
Neither $X,Y$ nor $f$ is required to be projective.
For proper pseudo-coherent morphisms of quasi-compact
quasi-separated schemes, $Rf_*$ preserves pseudo-coherent complexes.
If the morphism also has finite Tor dimension, it preserves perfect
complexes. This includes proper lci morphisms and nonnoetherian bases.

For complex analytic spaces, coherent sheaves likewise have a natural
transformation $\alpha_X:G_0^{\rm hol}(X)\to K^{\rm top}_0(X)$ to
topological $K$-homology, compatible with proper direct images.
On a complex manifold it sends a vector bundle to its Dolbeault
$K$-homology class. This includes possibly singular and nonprojective
complex spaces.
\end{theorem}
\begin{proof}
The algebraic transformation, its covariance and module formula are
\cite[Riemann--Roch Theorem (1)--(3), pp.~288--289]{BFultonGillet}.
Its Corollary 2 is exactly the displayed smooth formula; the virtual
class of $Rf_*E$ exists by perfection. Their argument uses projective
envelopes and the first stages of higher $K$-theory to descend the
projective Riemann--Roch transformation, so no unstated global
projective embedding is needed. Coherent sheaves represent bounded
pseudo-coherent complexes by alternating sums on these noetherian
schemes. The general preservation theorem is
\cite[Section 2, Kiehl's finiteness theorem on pp.~5--6;
Proposition 2.1 and Example 2.2(a)]{BLipmanNeeman2007}:
pseudo-coherence is preserved by proper pseudo-coherent maps,
and finite Tor dimension is preserved here by the projection formula.
The historical formulation in
\cite[Expos\'e III, Conjecture 2.1 and Theorem 2.2, pp.~236--237]{BSGA6}
proved the noetherian-base and projective cases and left the general
proper case conjectural. Thus the later finiteness input is substantive;
it is not obtained merely by reading ``SGA~6'' as a theorem of
unrestricted scope.
The independent derived-geometric proof for lci maps is
\cite[Theorems 0.3 and 2.1]{BToen2012}. Proper lci morphisms are
pseudo-coherent of finite Tor dimension. The proof descends to a
noetherian \emph{derived} model; an underived noetherian approximation
need not retain the lci property. This distinction removes a tempting
but invalid reduction in this nonnoetherian problem.

The analytic transformation and proper covariance are
\cite[Riemann--Roch Theorem (a)--(b), p.~149]{BLevy1987}; its actual
construction and functoriality proof occupy Section 3, pp.~168--186.
For smooth $X$, $K$-theoretic Poincar\'e duality interprets the bundle
as its Dolbeault class, accounting for the Todd class in ordinary
homology. In the singular case the target is $K$-\emph{homology},
not the vector-bundle ring $K^0(X)$; this makes proper functoriality
well typed. These are precise versions of the algebraic and analytic
Riemann--Roch programmes in subject 10; they do not impose a smooth
hypothesis on the singular transformation $\tau$ or $\alpha$.
\end{proof}

\begin{theorem}[Known with checked deduction; stack part of subject 14]\label{thm:B3-stacks}
Let $S$ be locally noetherian with G-ring local rings, and let
$\mathcal X$ be a category fibred in groupoids over $S$ satisfying the
usual set-size condition. Suppose its diagonal is representable by
algebraic spaces, it is an \'etale stack, and it commutes with filtered
limits of affine schemes. Suppose also that its deformation categories
satisfy the Rim--Schlessinger fibre-product condition and have finite
dimensional tangent and infinitesimal-automorphism spaces, every
formal object over a complete noetherian local $S$-algebra is effective,
and versality is open in finite type families. Then $\mathcal X$ is
an algebraic stack, locally of finite presentation over $S$.
\end{theorem}
\begin{proof}
This is the precise form of Artin's criterion in
\cite[Lemma 98.17.1, tag 07Y4, with axioms in Section 98.14,
tag 07XJ]{BStacks}. The Rim--Schlessinger condition means that the
functor on Artinian local algebras preserves the appropriate
fibre-product diagrams along surjections; effectiveness means that a
compatible system over $R/\mathfrak m^n$ comes from an object over $R$.
These are substantive hypotheses, not consequences of being a stack.
The proof first constructs versal formal deformations using the
finite-dimensional conditions, algebraizes by effectiveness, and
uses Artin approximation over the G-ring base to obtain finite type
schemes $U\to\mathcal X$. Openness of versality makes these maps
smooth near the chosen points. Taking one for each finite type
residue-field object gives a smooth surjective atlas: its image is
open and contains every finite type point of every test scheme.
The representable diagonal makes the atlas fibre products algebraic
spaces. The resulting atlas is locally of finite presentation, giving
the asserted finiteness over $S$. This settles a concrete
representability criterion and does not falsely claim that arbitrary
stacks are representable by schemes.
\end{proof}

\begin{theorem}[Known with checked deduction; artinian quotient part of subject 18]\label{thm:B3-artinian}
Let $R\rightrightarrows U=\operatorname{Spec}A$ be an equivalence
relation of schemes with both projections finite locally free.
Write $R=\operatorname{Spec}B$ and
\[
 C=\{a\in A:s^*a=t^*a\text{ in }B\}.
\]
Then the fppf quotient is the affine scheme $\operatorname{Spec}C$,
$U\to\operatorname{Spec}C$ is finite faithfully flat, and
$R=U\times_{\operatorname{Spec}C}U$.
In particular this applies to artinian $U$; its quotient is again
artinian. The same conclusion of schematic representability holds
for nonaffine $U$ if every equivalence class is contained in an affine
open, as is the case for $U$ quasi-projective over an affine scheme.
\end{theorem}
\begin{proof}
The first assertion is \cite[Proposition 39.23.9, tag 03BM]{BStacks},
whose proof identifies $A\otimes_C A\to B$ and establishes faithful
finite local freeness, not merely existence of an invariant subring.
It therefore proves representability of the \emph{sheaf} quotient.
When $A$ is artinian, descending chains of ideals of $C$ stabilize
after tensoring with the faithfully flat finite $C$-algebra $A$,
and hence stabilize already in $C$. Thus $C$ is artinian.
For the last assertion each point has a relation-invariant affine
neighbourhood by Lemma 39.24.1 (03JE). The affine quotients glue
along open subschemes, proving Proposition 66.14.1 (07S6).
Every scheme quasi-projective over an affine base has an affine
neighbourhood of any finite set of points, and the classes here are
finite. The word on
66/003 was rechecked on the official scan and is ``artinien''.
The statement is about equivalence relations, hence free actions
when expressed by a group action; a nonfree action has stabilizers
and must not be substituted into the assertion $R=U\times_MU$.
\end{proof}

\begin{theorem}[Known with checked deduction; coherent part of subject 6]\label{thm:B3-holomorphic-lefschetz}
Let $X$ be a compact complex manifold, $E$ a holomorphic vector
bundle, $f:X\to X$ holomorphic, and $u:f^*E\to E$ a holomorphic
bundle map. Suppose every fixed point $x$ is nondegenerate:
$\det_\C(1-df_x)\ne0$. Then the induced maps on coherent cohomology
satisfy
\[
 \sum_q(-1)^q\operatorname{Tr}(u\circ f^*:H^q(X,E))
 =\sum_{f(x)=x}\frac{\operatorname{Tr}(u_x)}
                         {\det_\C(1-df_x)}.
\]
More generally an elliptic complex on a compact smooth manifold,
with a geometric endomorphism over a map having only nondegenerate
fixed points, has local contribution
$\sum_q(-1)^q\operatorname{Tr}(u_{q,x})/
 |\det_\R(1-df_x)|$.
\end{theorem}
\begin{proof}
The latter formula is \cite[Theorem A, pp.~377--378]{BAtiyahBott1967};
the full paper proves the analytic localization, rather than merely
announcing it. Apply it to the Dolbeault resolution
$E\otimes\Omega_X^{0,\bullet}$, whose ellipticity and Dolbeault
cohomology identify the Lefschetz number with the left side above.
At a fixed point the alternating trace on the antiholomorphic
exterior algebra is $\det_\C(1-\overline{df_x})$. The real determinant
in the denominator is
$|\det_\C(1-df_x)|^2$, so cancellation gives precisely the displayed
holomorphic denominator. Nondegeneracy makes the fixed points
isolated; compactness makes their number finite. The theorem therefore
gives an explicit coherent-coefficient formula without an algebraic
or projectivity assumption. Positive-dimensional or degenerate
fixed loci require the more general local-term formalism and are
not described by this isolated-point denominator. The discrete
\'etale counterpart was stated in Theorem~\ref{thm:B3-trace}.
\end{proof}

\begin{theorem}[Known with checked deduction; subject 15]\label{thm:B3-homotopy}
Every $\infty$-topos $\mathcal X$ has a shape, the pro-space
corepresenting the accessible left exact functor
$K\mapsto\Gamma_{\mathcal X}(\Delta K)$.
A pointed object $A\in\mathcal X$ has internal homotopy sheaves
$\pi_n(A)=\tau_{\leq0}\Omega^nA$; fibre sequences give the usual
long exact sequence, with the group action interpretation at
$\pi_0$. If $A$ is $(n-1)$-connected and $n\geq1$, then
\[
 \pi_n(A)^{\mathrm{ab}}\otimes M
       \simeq\widetilde H_n(A;M)
\]
for every internal abelian group $M$.
For a finite type $\C$-scheme $X$, the profinite completion of its
\'etale shape is the profinite completion of the singular homotopy
type of $X^{\rm an}$.

Infinitesimal homotopy sheaves are possible on the fppf site:
in characteristic $p$, $B\alpha_p$ has $\pi_1=\alpha_p$.
They are not detected by the small \emph{\'etale} site of a reduced
scheme. The distinction between internal homotopy sheaves and
external pro-homotopy groups of the shape is essential.
\end{theorem}
\begin{proof}
The shape and internal homotopy constructions are
\cite[Definitions 7.1.6.1 and 7.1.6.3; Definition 6.5.1.1 and
Remarks 6.5.1.3--6.5.1.6]{BLurieHTT}. The exact sequence is constructed
by the decomposition of a sphere into two hemispheres; exactness
reduces through a left exact localization to the objectwise statement
for spaces (Remark 6.5.1.5). The internal Hurewicz theorem is
\cite[Theorem H and Theorem 3.12]{BChristensenScoccola}; its statement
and proof apply in every $\infty$-topos, with the abelianization
needed only for $n=1$.

For the complex comparison, finite covering spaces agree by
\cite[Expos\'e XII, Theorem 5.1, original pp.~332--333]{BSGA1},
and the cohomology of finite local systems agrees by
\cite[Expos\'e XVI, Theorem 4.1]{BSGA4XVI}. Induct on the finite
Postnikov tower of a finite homotopy type: the covering-space
statement gives its $1$-truncation, and the obstruction and lift
spaces for each following Eilenberg--MacLane fibre are the compared
cohomology groups with finite local coefficients. This identifies
mapping spaces into every truncated space with finite homotopy
groups, which is exactly the defining test for profinite shape
equivalence. The uncompleted statement is not inferred. In
particular Lurie's Example 7.1.6.9 distinguishes the original
Artin--Mazur pro-homotopy category construction from the shape,
its hypercompletion and its protruncation.

For the infinitesimal example, the fppf exact sequence
$0\to\alpha_p\to\mathbf G_a\xrightarrow{F}\mathbf G_a\to0$
gives, on $X=\operatorname{Spec}k[t]$ with $k$ algebraically closed,
\[
 H^1_{\mathrm{fppf}}(X,\alpha_p)=k[t]/k[t]^p\ne0.
\]
Surjectivity of $F$ as a sheaf follows from the finite free cover
adjoining a $p$th root; $H^1(X,\mathbf G_a)=0$ is affine
quasi-coherent acyclicity. The class of $t$ is nonzero because a
polynomial $p$th power has only exponents divisible by $p$.
But every scheme \'etale over this reduced $X$ is reduced, so its
only function with $a^p=0$ is zero: $\alpha_p$ restricts to the zero
sheaf on the small \'etale site. The classifying object construction
gives $\Omega B\alpha_p=\alpha_p$ and the claimed homotopy sheaf.
This explicitly supplies an infinitesimal kernel while specifying
the topology in which it exists.
\end{proof}

\begin{theorem}[Known with checked deduction; subject 23]\label{thm:B3-chow-pullback}
Let $Y$ be a regular separated scheme of finite type over $\Z$.
For each morphism $f:X\to Y$ of separated finite type $\Z$-schemes
there is a canonical rational bivariant orientation $[f]$, uniquely
characterized by $[f]([Y])=[X]$. For irreducible $X,Y$ it has degree
$-\dim_YX$ and gives, in the relative dimension grading,
\[
 f^*:\mathrm{CH}_k(Y)_\Q\longrightarrow
       \mathrm{CH}_{k+\dim_\Z X-\dim_\Z Y}(X)_\Q.
\]
It is compatible with proper pushforward, flat pullback and refined
Gysin maps; compositions of orientations agree whenever the
intermediate target is also regular. In particular $X$ need not be
regular, and $f$ need not be flat or lci. The coefficient field
$\Q$ and the finite type hypotheses are retained.
\end{theorem}
\begin{proof}
The precise meaning of orientation is
\cite[Definition 4.1.12 and Lemma 4.1.13, pp.~29--30]{BWeber2015}.
Theorem 4.2.1, pp.~37--38, proves that a regular target is orienting
whenever regular projective alterations of invertible degree exist
up to its dimension; Corollary 4.2.3, p.~39, gives the unconditional
assertion over $\operatorname{Spec}\Z$ with rational coefficients.
Here is the construction underlying this result. It suffices by
Proposition 4.1.16 to orient projective $g:V\to Y$ with $V$ integral.
Flatten $V$ over a blowup of its image, and take a regular projective
alteration of that blowup of degree $d$. In the resulting diagram
there is a flat map $g':V'\to W$, a projective lci map $h:W\to Y$
between regular schemes, and a proper generically finite
$h':V'\to V$ of degree $d$. The bivariant class
\[
             [g]=d^{-1}h'_*([g']\cdot[h])
\]
sends $[Y]$ to $[V]$. Bivariant descent and the orienting criterion
prove its uniqueness, hence independence of the chosen flattening
and alteration and compatibility with composition. Thus this is a
well-defined operation on Chow groups, not a proposed inverse image
of cycle representatives whose rational-equivalence compatibility
has been left unchecked. In the smooth-over-a-field case, the graph
construction and refined intersection already give integral
pullback. The displayed arithmetic result makes explicit the
rational replacement when the sole smoothness condition is
regularity of the target. No assertion of integral pullback for
arbitrary regular bases is silently added.
\end{proof}

\subsection{B.3: the nine headings in the \emph{Pour m\'emoire} diagram}

In citations to \cite{BILO2014}, chapter and theorem numbers are
the primary locators. Page numbers refer to the complete opened
author version of 16 July 2012, whose pagination differs from the
published 2014 volume. The bibliography records both versions
explicitly.

\begin{theorem}[Known with checked deduction; resolution and its alteration replacement]\label{thm:B3-resolution}
Every generically reduced noetherian quasi-excellent $\Q$-scheme
admits a finite sequence of blowups with regular centres, disjoint
from the inverse image of its regular locus, whose final scheme is
regular. The construction is functorial for regular morphisms, up to
omitting empty blowups.
For a separated finite type scheme $X$ over any field $k$, a nowhere
dense closed $Z\subset X$, and a prime $\ell\ne\operatorname{char}k$,
there are a finite extension $k'/k$ of degree prime to $\ell$ and a
proper surjective generically finite alteration $X'\to X$ of degree
prime to $\ell$, with $X'$ smooth quasi-projective over $k'$ and
inverse image of $Z$ the support of a strict normal-crossings divisor.
The alteration assertion does not assert a birational resolution in
arbitrary characteristic and dimension.
\end{theorem}
\begin{proof}
The first is \cite[Theorem 1.2.1]{BTemkin2012}, including all four
conditions; its proof globalizes functorial characteristic-zero
resolution through formal desingularization and algebraization.
The second is Gabber's refinement of alteration theory,
\cite[Introduction, Theorem 3(1), and Expos\'e X, Theorem 2.1]{BILO2014}.
The degree condition permits trace arguments with $\ell$-primary
coefficients; it does not turn the alteration into a modification.
These are the precise general results used for the cohomological
arrows in the diagram. Replacing the characteristic-zero hypothesis
in the first statement by arbitrary characteristic is not a conclusion
of either theorem and is not claimed in this report.
\end{proof}

\begin{theorem}[Known with checked deduction; absolute cohomological purity]\label{thm:B3-purity}
Let $i:Z\hookrightarrow X$ be a closed immersion of constant
codimension $c$ between regular noetherian schemes, and let
$\Lambda=\Z/n\Z$ with $n$ invertible on $X$. The cycle-class Gysin
morphism is an isomorphism
\[
                 Ri^!\Lambda\simeq\Lambda(-c)[-2c].
\]
In particular $\mathcal H^q_Z(\Lambda)$ vanishes for $q\ne2c$;
in degree $2c$ it is $i_*\Lambda(-c)$.
\end{theorem}
\begin{proof}
This is Gabber's absolute-purity theorem, with complete proof in
\cite[Expos\'e XVI, Theorem 3.1.1, p.~242]{BILO2014}. The stated
normalization follows by solving its formula
$\Lambda\simeq i^!\Lambda(c)[2c]$ for $i^!\Lambda$.
Applying the exact direct image $i_*$ identifies local cohomology
with $i_*Ri^!$, giving the vanishing. Regularity of both schemes
and invertibility of $n$ are part of the assertion; a singular
closed support is not assigned the same purity isomorphism.
\end{proof}

\begin{theorem}[Known with checked deduction; regular base change]\label{thm:B3-regular-bc}
In a Cartesian square with $f:T\to S$ quasi-compact and
quasi-separated and $g:S'\to S$ a regular morphism of noetherian
schemes, let $f':T'\to S'$ and $g':T'\to T$ be the projections.
For every \'etale $\Z/n\Z$-sheaf $F$ on $T$, where $n$ is invertible
on $S$, the base-change map is an isomorphism
\[
                    g^*Rf_*F\simeq Rf'_*g'^*F.
\]
\end{theorem}
\begin{proof}
The full statement and proof are
\cite[Expos\'e XIV, Lemma 2.5.3, pp.~216--217]{BILO2014}.
Zariski locally Popescu's theorem writes the regular map as a
filtered limit of smooth maps with affine transition maps.
Smooth base change gives the assertion at each finite stage,
and continuity of \'etale cohomology for quasi-compact
quasi-separated maps passes it to the limit. No finite type
assumption on the regular map is introduced. For schemes with
complete excellent local rings, the completion map is regular,
so the theorem also covers the regular-completion reductions
used later in that volume.
\end{proof}

\begin{theorem}[Known with checked deduction; local duality and biduality]\label{thm:B3-local-biduality}
Let $X$ be a noetherian excellent $\Z[1/n]$-scheme with a dimension
function $\delta$, and put $\Lambda=\Z/n\Z$. There is a normalized
potential dualizing complex $K_{X,\delta}$, unique with its
normalization, belonging to $D^b_{\mathrm{ctf}}(X_{\mathrm{\acute et}},\Lambda)$.
The functor $D_K=R\mathcal Hom_\Lambda(-,K_{X,\delta})$ preserves
$D^b_c$, and for every $M\in D^b_c$ the canonical map
\[
                            M\longrightarrow D_KD_KM
\]
is an isomorphism. The complex has finite quasi-injective dimension
if and only if $X$ has finite Krull dimension; in that case it is
also a dualizing complex in the SGA~5 sense.
\end{theorem}
\begin{proof}
Existence and normalized uniqueness are
\cite[Expos\'e XVII, Theorem 5.1.1, p.~287]{BILO2014}; all four
assertions concerning bounded constructibility, Tor dimension,
quasi-injective dimension and biduality are Theorem 6.1.1,
pp.~293--294. The normalization prescribes the local cohomology
at each geometric point by its dimension function. Absolute
purity supplies the regular case; the proof then glues potential
dualizing complexes and uses constructibility for open direct
images to extend biduality across closed strata. Thus the diagram's
``dualit\'e locale = bidualit\'e'' is an actual involution theorem
on the specified constructible derived category, without an
unstated smoothness assumption on $X$.
\end{proof}

\begin{theorem}[Known with checked deduction; finite type finiteness]\label{thm:B3-finiteness}
Let $X$ be noetherian quasi-excellent, $f:Y\to X$ of finite type,
and $F$ a constructible \'etale $\Z/n\Z$-sheaf on $Y$, with $n$
invertible on $X$. Then every $R^qf_*F$ is constructible and
$R^qf_*F=0$ for all sufficiently large $q$. Equivalently
$Rf_*$ preserves bounded constructible complexes. No properness
assumption on $f$ or finite Krull dimension assumption on $X$
is needed in this theorem.
\end{theorem}
\begin{proof}
This is exactly \cite[Expos\'e XIII, Theorem 1.1.1 and 1.1.2,
p.~195]{BILO2014}. The source separates the two assertions:
Section 2 proves constructibility using weak local uniformization,
absolute purity, generic constructibility and cohomological
descent; Section 3 obtains boundedness in general using
prime-to-$\ell$ alterations. Merely quoting proper base change
would establish neither the stated nonproper generality nor this
boundedness. The result directly settles the diagram's finite type
finiteness heading for torsion coefficients invertible on the base.
The earlier original theorem
\cite[Th\'eor\`emes de finitude, Theorem 1.1 and Corollary 1.5,
pp.~233--234]{BSGA4Half}, now visually opened, assumes a regular
noetherian base of dimension zero or one. It is not substituted for
the stronger quasi-excellent-base theorem used here.
\end{proof}

\begin{theorem}[Known with checked deduction; affine cohomological dimension]\label{thm:B3-affine-cd}
Let $X$ be noetherian quasi-excellent with dimension function
$\delta$, let $f:Y\to X$ be affine of finite type, and let $F$ be
a constructible \'etale $\Z/n\Z$-sheaf with $n$ invertible on $X$.
Set $\delta_Y(y)=\delta(f(y))+\operatorname{trdeg}_{\kappa(f(y))}\kappa(y)$
and $\delta(F)=\sup\{\delta_Y(y):F_{\bar y}\ne0\}$.
Then
\[
 \sup\{\delta(x):(R^qf_*F)_{\bar x}\ne0\}\leq\delta(F)-q.
\]
For an affine variety of dimension $d$ over a separably closed field,
this implies $H^q(Y,F)=0$ for $q>d$.
\end{theorem}
\begin{proof}
The precise relative statement is
\cite[Expos\'e XV, Theorem 1.1.2 and Corollary 1.2.2, p.~219]{BILO2014}.
With $X$ a point and $\delta=0$, the support dimension is at most
$d$, giving the final assertion. The proof reduces to an affine
open of a strictly local quasi-excellent scheme, passes to its
completion by regular base change, and uses the affine Lefschetz
argument with alterations. The use of a dimension function in the
relative assertion matters: the absolute bound cannot simply be
replaced by a bound on the dimensions of the fibres of an arbitrary
affine morphism.
\end{proof}

\begin{theorem}[Known with checked deduction; Grauert finiteness]\label{thm:B3-grauert}
For a proper holomorphic map $f:X\to Y$ of complex analytic
spaces and a coherent analytic $\mathcal O_X$-module $F$,
every $R^qf_*F$ is coherent. In particular cohomology of a
coherent sheaf on a compact complex space is finite dimensional.
\end{theorem}
\begin{proof}
This is \cite[Section 6, Hauptsatz I, p.~59]{BGrauert1960}, read
with the corrections in \cite[pp.~35--36]{BGrauert1963}.
The proof establishes finite local generators and finite relations
for direct-image cohomology by estimates on complexes over
polycylinders; coherence is local on $Y$, so an embedding of a
neighbourhood in a polydisc gives the general analytic-space case.
When $Y$ is a point, coherent modules are finite dimensional complex
vector spaces. The correction replaces an insufficient uniform
radius condition by a triangular one and repairs the corresponding
induction. It also warns that a locus asserted closed in
Section 7, Satz 2 need not be closed. No such closed-locus assertion
is used here, and the corrected proper mapping theorem is the one
being applied.
\end{proof}

\begin{theorem}[Known with checked deduction; rigid analytic theory]\label{thm:B3-rigid}
Let $K$ be complete for a nontrivial rank-one nonarchimedean
valuation. On an affinoid rigid space $X=\operatorname{Sp}A$,
finite $A$-modules correspond to coherent analytic sheaves and
have no higher cohomology. For every proper map of rigid
$K$-spaces, all higher direct images of coherent sheaves are
coherent. If $R$ is the valuation ring, taking generic fibre
induces an equivalence between quasi-paracompact admissible
formal $R$-schemes localized at admissible formal blowups and
quasi-paracompact quasi-separated rigid $K$-spaces.
\end{theorem}
\begin{proof}
Exact locators in the full opened text are
\cite[Section 4.3, Theorem 10 and Corollary 11, pp.~90--91;
Section 6.1, Theorem 4, p.~119;
Section 6.3, Theorem 9, p.~132;
Section 8.4, Theorem 3, p.~204]{BBosch2014}.
Tate's acyclicity makes affinoid functions and finite modules
sheaves for the admissible topology; Kiehl's theorem identifies
all coherent sheaves on an affinoid with finite modules.
Kiehl's proper mapping theorem then supplies the rigid counterpart
of Grauert finiteness. Raynaud's theorem is a categorical
comparison, not just an association of a generic fibre: its proof
shows that every rigid space in the stated class has a formal
model, and that every rigid morphism lifts after an admissible
blowup of its source. The quasi-separatedness and
quasi-paracompactness conditions in the displayed equivalence
are retained.
\end{proof}

\begin{theorem}[Proved counterexample and known replacement; questions in SGA 1962]\label{thm:B3-sga2}
The reference is to SGA~2, and its proposed finiteness statements
cannot all be recorded as affirmative theorems. In particular
SGA~2, Expos\'e XIII, Conjecture 1.1 is false: a local cohomology
module of a finite module over a regular noetherian ring need not
have finitely generated $I$-annihilator.
A precise positive local Lefschetz conclusion is that every
noetherian local complete-intersection ring of dimension at least
four is parafactorial. Consequently a noetherian local
complete-intersection ring that is factorial at every prime of
height at most three is factorial.
\end{theorem}
\begin{proof}
For the negative assertion let $A=k[u,v,x,y]$,
$M=R=A/(ux-vy)$ and $I=(u,v)$. The \v{C}ech complex computes
\[
              H^2_I(M)=R_{uv}/(R_u+R_v).
\]
For every $n\geq0$ let $e_n$ be the class of
$x^n/(uv^{n+1})$. Multiplication by $u$ is zero modulo $R_v$;
multiplication by $v$ is zero modulo $R_u$, since
$x^n/(uv^n)=y^n/u^{n+1}$. Likewise $xe_n$ is represented by
$y^{n+1}/u^{n+2}$ and $ye_n$ by $x^{n+1}/v^{n+2}$, so every
$e_n$ is annihilated by $(u,v,x,y)$.

To verify nonvanishing and independence, embed the toric domain
$R$ into $k[a,b,c,d]$ by
$u=ac$, $v=ad$, $x=bd$, $y=bc$.
Its only defining relation is $ux-vy$ (the kernel is a prime of
height one containing this irreducible polynomial).
In the resulting Laurent monomial algebra, $e_n$ is
$a^{-n-2}b^nc^{-1}d^{-1}$. Every monomial in $R_u$ has
nonnegative $d$ exponent and every monomial in $R_v$ has
nonnegative $c$ exponent. Distinct displayed monomials are
therefore linearly independent modulo $R_u+R_v$.
Thus $\operatorname{Hom}_A(A/I,H^2_I(M))$ contains an infinite
dimensional subspace annihilated by the maximal ideal
$(u,v,x,y)$. If this Hom module were finitely generated, its
maximal-ideal annihilator would be a finite submodule over the
noetherian ring $A$, hence finite dimensional over $k$, a
contradiction. This explicitly refutes the conjecture as stated
in \cite[Expos\'e XIII, Conjecture 1.1, pp.~173--174 of the original;
editorial discussion and Hartshorne attribution]{BSGA2}.
The original counterexample is
\cite[Section 3, pp.~149--151]{BHartshorne1970}, with
$A=k[x,y][[u,v]]$ and $M=A/(xu+yv)$; those pages have been
visually checked. The polynomial-ring variant above supplies
its own proof without a completion comparison being assumed.

For the positive assertion use
\cite[Expos\'e XI, Theorem 3.13(ii) and Corollary 3.14,
original p.~132]{BSGA2}. Parafactoriality here includes depth
at least two and $\operatorname{Pic}(\operatorname{Spec}A\setminus
\{\mathfrak m\})=0$. To obtain factoriality, induct on dimension:
the hypothesis and localization of a complete intersection make
all punctured local rings factorial by induction, and
parafactoriality removes the remaining divisor class at the
closed point. This is the proof of Corollary 3.14. These statements
address concrete questions within the SGA~2 programme; the
undifferentiated heading is not a claim that every conjecture in
that volume has a positive answer.
\end{proof}
